\section{Entropy functionals and their monotonicity}
\dcref{I.1--I.6}

In Sections I.1.1--I.5, $M$ is closed. We use the Einstein summation
convention and $\triangle=\Div\grad$. Let ${\mathcal M}$ be the space of
smooth metrics on $M$, with $T_g{\mathcal M}$ the symmetric covariant
$2$-tensors, and let $dV$ be the volume density.

\subsection{The ${\mathcal F}$-functional and its monotonicity} \label{I.1.1}
\dcref{I.1.1}
\label{Derivation of first equation of Section 1.1}

\begin{definition}[F-entropy definition]
  \label{def:kl-f-entropy-f-entropy-definition}
  \source{kleiner-lott}
  \dcref{Definition 5.1}
  \group{Kleiner--Lott Entropy Functionals}
The ${\mathcal F}$-functional is
\begin{equation}
{\mathcal F}(g,f)=\int_M(R+|\nabla f|^2)e^{-f}\,dV.
\end{equation}
\end{definition}

For $v_{ij}\in T_g{\mathcal M}$ and $h\in C^\infty(M)$, put
$v=g^{ij}v_{ij}$. Then
\begin{proposition}[first variation of F-entropy]
  \label{prop:kl-f-entropy-first-variation-f-entropy}
  \source{kleiner-lott}
  \dcref{Proposition 5.3}
  \group{Kleiner--Lott Entropy Functionals}
\begin{equation}
\delta {\mathcal F}(v_{ij},h)=\int_Me^{-f}\left[-v_{ij}(R_{ij}+\nabla_i\nabla_jf)
+\left({v\over2}-h\right)(2\triangle f-|\nabla f|^2+R)\right]dV.
\end{equation}
\end{proposition}
\begin{proof}
\sketch
\begin{equation}
\delta R=-\triangle v+\nabla_i\nabla_jv_{ij}-R_{ij}v_{ij},\qquad
\delta|\nabla f|^2=-v^{ij}\nabla_if\nabla_jf+2\langle\nabla f,\nabla h\rangle,
\end{equation}
\begin{equation}\label{measvar}
\delta(e^{-f}dV)=\left({v\over2}-h\right)e^{-f}dV.
\end{equation}
Using integration by parts and
\begin{equation}\label{id}
\triangle e^{-f}=(|\nabla f|^2-\triangle f)e^{-f},
\end{equation}
gives
\begin{equation}\label{Fvariation}
\delta {\mathcal F}=\int_Me^{-f}\left[-v_{ij}(R_{ij}+\nabla_i\nabla_jf)
+\left({v\over2}-h\right)(2\triangle f-|\nabla f|^2+R)\right]dV.
\end{equation}
\end{proof}

\begin{equation}\label{putting}
\delta{\mathcal F}=\int_Me^{-f}\left[-\triangle v+\nabla_i\nabla_jv_{ij}-R_{ij}v_{ij}
-v_{ij}\nabla_if\nabla_jf+2\langle\nabla f,\nabla h\rangle
+(R+|\nabla f|^2)\left({v\over2}-h\right)\right]dV.
\end{equation}

Fix a smooth measure $dm$ and impose $e^{-f}dV=dm$. Then
\begin{equation}
d{\mathcal F}^m(v_{ij})=-\int_Me^{-f}v_{ij}(R_{ij}+\nabla_i\nabla_jf)dV.
\end{equation}
With
\begin{equation}
\langle v_{ij},v_{ij}\rangle_g={1\over2}\int_Mv^{ij}v_{ij}\,dm,
\end{equation}
the gradient flow and the induced $f$-equation are
\begin{equation}\label{Fflow}
(g_{ij})_t=-2(R_{ij}+\nabla_i\nabla_jf),
\end{equation}
\begin{equation}
f_t={\partial\over\partial t}\log{dV\over dm}=-\triangle f-R.
\end{equation}
Moreover,
\begin{equation}\label{Fchange}
{d\over dt}{\mathcal F}^m=2\int_M|R_{ij}+\nabla_i\nabla_jf|^2\,dm.
\end{equation}
After the time-dependent diffeomorphism generated by $V=\nabla f$,
\begin{align}\label{newevolution}
(g_{ij})_t&=-2R_{ij},\\
f_t&=-\triangle f-R+|\nabla f|^2.\notag
\end{align}
The monotonicity formula becomes
\begin{equation}\label{Fder}
{d\over dt}{\mathcal F}(g(t),f(t))=2\int_M|R_{ij}+\nabla_i\nabla_jf|^2e^{-f}dV,
\end{equation}
and $\int_Me^{-f}dV$ is constant. Equivalently,
\begin{equation}\label{backwards}
{\partial\over\partial t}e^{-f}=-\triangle e^{-f}+Re^{-f}.
\end{equation}
Thus, given Ricci flow on $[t_1,t_2]$, solve (\ref{backwards}) backward from
$f(t_2)$ to obtain a pair satisfying (\ref{Fder}).

\subsection{Basic example for the ${\mathcal F}$-functional}

On static Euclidean $\mathbb R^n$, let $\tau=t_0-t>0$ and
\begin{equation}
f(t,x)={|x|^2\over4\tau}+{n\over2}\log(4\pi\tau),qquad
e^{-f}=(4\pi\tau)^{-n/2}e^{-|x|^2/(4\tau)}.
\end{equation}
Then $(g,f)$ satisfies (\ref{newevolution}), and
\begin{equation}\label{Gaussianintegral}
\int_{\mathbb R^n}e^{-|x|^2/(4\tau)}dV=(4\pi\tau)^{n/2},
\end{equation}
\begin{equation}\label{Gaussianintegral2}
\int_{\mathbb R^n}{|x|^2\over4\tau^2}e^{-|x|^2/(4\tau)}dV
=(4\pi\tau)^{n/2}{n\over2\tau}.
\end{equation}
Consequently ${\mathcal F}(t)=n/(2\tau)=n/[2(t_0-t)]$.

\subsection{The $\lambda$-invariant and its applications}
\label{I.2.2}
\dcref{I.2.2}

\begin{proposition}[uniqueness for F-entropy]
  \label{prop:kl-invariant-its-applications-uniqueness-f-entropy}
  \source{kleiner-lott}
  \dcref{Proposition 7.1}
  \group{Kleiner--Lott Entropy Functionals}
For fixed $g$, there is a unique minimizer $\bar f$ of ${\mathcal F}(g,f)$
subject to $\int_Me^{-f}dV=1$.
\end{proposition}
\begin{proof}
\sketch
With $\Phi=e^{-f/2}$,
\begin{equation}
{\mathcal F}=\int_M(4|\nabla\Phi|^2+R\Phi^2)dV
=\int_M\Phi(-4\triangle\Phi+R\Phi)dV,qquad\int_M\Phi^2dV=1.
\end{equation}
Thus the minimum is the lowest eigenvalue of $-4\triangle+R$, whose normalized
positive eigenfunction is unique \cite[Chapter XIII.12]{Reed-Simon}.
\end{proof}

\begin{definition}[lambda invariant]
  \label{def:kl-invariant-its-applications-lambda-invariant}
  \source{kleiner-lott}
  \dcref{Definition 7.4}
  \group{Kleiner--Lott Entropy Functionals}
$\lambda(g)={\mathcal F}(g,\bar f)$.
\end{definition}

\begin{proposition}[monotonicity for entropy]
  \label{prop:kl-invariant-its-applications-monotonicity-entropy}
  \source{kleiner-lott}
  \dcref{Proposition 7.5}
  \group{Kleiner--Lott Entropy Functionals}
  (cf. I.2.2) \label{mono}
Along Ricci flow, $\lambda(g(t))$ is nondecreasing.
\end{proposition}
\begin{proof}
\sketch
Given $t_1<t_2$, solve (\ref{backwards}) backward with
$u(t_2)=e^{-\bar f(t_2)}$ and put $u=e^{-f}$. Positivity follows by duality
with the forward heat kernel, while $\int_Mu(t)dV(t)=1$. Formula (\ref{Fder})
and the minimizing property give
\[\lambda(t_1)\le {\mathcal F}(g(t_1),f(t_1))le {\mathcal F}(g(t_2),f(t_2))=\lambda(t_2).\]
\end{proof}

\begin{definition}[steady breather definition]
  \label{def:kl-invariant-its-applications-steady-breather-definition}
  \source{kleiner-lott}
  \dcref{Definition 7.9}
  \group{Kleiner--Lott Entropy Functionals}
A steady breather is a Ricci flow on $[t_1,t_2]$ with
$g(t_2)=\phi^*g(t_1)$ for some $\phi\in\Diff(M)$.
\end{definition}

\begin{proposition}[rigidity of steady breathers]
  \label{prop:kl-invariant-its-applications-rigidity-steady-breathers}
  \source{kleiner-lott}
  \dcref{Proposition 7.10}
  \group{Kleiner--Lott Entropy Functionals}
  (cf. I.2.2) \label{steadybreather}
A steady breather is a gradient steady soliton.
\end{proposition}
\begin{proof}
  \uses{prop:kl-invariant-its-applications-monotonicity-entropy}
\sketch
Diffeomorphism invariance gives $\lambda(t_2)=\lambda(t_1)$, hence equality in
the preceding monotonicity proof. Equation (\ref{Fder}) yields
$R_{ij}+\nabla_i\nabla_jf=0$, and (\ref{newevolution}) is therefore the steady
soliton equation.
\end{proof}

\begin{lemma}[differential inequality for entropy]
  \label{lem:kl-invariant-its-applications-differential-inequality-entropy}
  \source{kleiner-lott}
  \dcref{Lemma 7.11}
  \group{Kleiner--Lott Entropy Functionals}
  (cf. Proposition I.1.2) \label{2/3}
\begin{equation}
{d\lambda\over dt}\ge {2\over n}\lambda^2(t).
\end{equation}
\end{lemma}
\begin{proof}
  \uses{prop:kl-invariant-its-applications-monotonicity-entropy}
\sketch
The proof above gives
\begin{equation}\label{dlambdadt}
{d\lambda\over dt}\ge2\int_M|R_{ij}+\nabla_i\nabla_j\bar f|^2e^{-\bar f}dV
\ge {2\over n}\int_M(R+\triangle\bar f)^2e^{-\bar f}dV.
\end{equation}
Since $\int e^{-\bar f}dV=1$, Cauchy--Schwarz and (\ref{id}) give
\begin{equation}\label{using}
\int_M(R+\triangle\bar f)e^{-\bar f}dV
=\int_M(R+|\nabla\bar f|^2)e^{-\bar f}dV=\lambda(t),
\end{equation}
which proves the claim.
\end{proof}

\subsection{The rescaled $\lambda$-invariant}
\dcref{I.2.3}

Put $\bar\lambda(g)=\lambda(g)V(g)^{2/n}$. Then
\begin{proposition}[monotonicity of the rescaled lambda invariant]
  \label{prop:kl-rescaled-invariant-monotonicity-rescaled-lambda-invariant}
  \source{kleiner-lott}
  \dcref{Proposition 8.1}
  \group{Kleiner--Lott Entropy Functionals}
  \label{expandingprop} (cf. Claim of I.2.3)
If $\bar\lambda(g(t))\le0$ at some time, then $d\bar\lambda/dt\ge0$.
\end{proposition}
\begin{proof}
\sketch
\begin{equation}
{d\bar\lambda\over dt}=V^{2/n}{d\lambda\over dt}-{2\over n}V^{(2-n)/n}\lambda\int_MR,dV.
\end{equation}
Testing ${\mathcal F}$ with the constant function gives
\begin{equation}\label{ttest}\lambda(t)\le V(t)^{-1}\int_MR,dV.
\end{equation}
For $\lambda\le0$, use (\ref{dlambdadt}), (\ref{ttest}), and
\[|S|^2\ge|S-({\tr S/n})g|^2+{1\over n}(\tr S)^2\]
to obtain
\begin{align}\label{cs}
{d\bar\lambda\over dt}\ge2V^{2/n}\bigg[&\int_M\left|R_{ij}+\nabla_i\nabla_j\bar f-{R+\triangle\bar f\over n}g_{ij}\right|^2e^{-\bar f}dV\\
&+{1\over n}\int_M(R+\triangle\bar f)^2e^{-\bar f}dV-{1\over n}\left(\int_M(R+\triangle\bar f)e^{-\bar f}dV\right)^2\bigg]\ge0.
\end{align}
\end{proof}

\begin{corollary}[rigidity at constant rescaled lambda]
  \label{cor:kl-rescaled-invariant-rigidity-at-constant-rescaled-lambda}
  \source{kleiner-lott}
  \dcref{Corollary 8.9}
  \group{Kleiner--Lott Entropy Functionals}
  \label{expp}
If $\bar\lambda$ is constant and nonpositive on $[t_1,t_2]$, the flow is a
gradient soliton.
\end{corollary}
\begin{proof}
\sketch
Equality in (\ref{cs}) gives $R+\triangle\bar f=\alpha(t)$ and
$R_{ij}+\nabla_i\nabla_j\bar f=\alpha(t)g_{ij}/n$, hence the flow evolves by
diffeomorphisms and dilations.
\end{proof}

\begin{definition}[expanding breather definition]
  \label{def:kl-rescaled-invariant-expanding-breather-definition}
  \source{kleiner-lott}
  \dcref{Definition 8.10}
  \group{Kleiner--Lott Entropy Functionals}
An expanding breather is a Ricci flow on $[t_1,t_2]$ with
$g(t_2)=c\phi^*g(t_1)$ for $c>1$ and $\phi\in\Diff(M)$.
\end{definition}

\begin{proposition}[rigidity of expanding breathers]
  \label{prop:kl-rescaled-invariant-rigidity-expanding-breathers}
  \source{kleiner-lott}
  \dcref{Proposition 8.11}
  \group{Kleiner--Lott Entropy Functionals}
An expanding breather is a gradient expanding soliton.
\end{proposition}
\begin{proof}
  \uses{cor:kl-rescaled-invariant-rigidity-at-constant-rescaled-lambda,prop:kl-rescaled-invariant-monotonicity-rescaled-lambda-invariant}
\sketch
Scale and diffeomorphism invariance give $\bar\lambda(t_2)=\bar\lambda(t_1)$.
Since $V(t_2)>V(t_1)$, (\ref{ttest}) implies $\bar\lambda<0$ somewhere; monotonicity
then makes it a negative constant, so apply Corollary \ref{expp}.
\end{proof}

\subsection{Gradient steady solitons on compact manifolds}
\dcref{I.2.4}

\begin{proposition}[compactness for Ricci solitons]
  \label{prop:kl-gradient-steady-solitons-compact-manifolds-compactness-ricci-solitons}
  \source{kleiner-lott}
  \dcref{Proposition 9.1}
  \group{Kleiner--Lott Entropy Functionals}
A gradient steady solution on a compact manifold is Ricci flat.
\end{proposition}
\begin{proof}
  \uses{prop:kl-invariant-its-applications-monotonicity-entropy}
\sketch
In the equality case, $e^{-f/2}$ is the normalized lowest eigenfunction:
\[(-4\triangle+R)e^{-f/2}=\lambda e^{-f/2},\qquad2\triangle f-|\nabla f|^2+R=\lambda.\]
Since $R+\triangle f=0$, $\triangle e^{-f}=-\lambda e^{-f}$. Integration gives
$\lambda=0$, and then $\int_M|\nabla e^{-f}|^2dV=0$; hence $f$ is constant and
$\Ric=0$.
\end{proof}

\subsection{Ricci flow as a gradient flow}

For $\Phi=e^{-\bar f/2}$, define
\begin{equation}\label{riemmetric}
\langle v_{ij},v_{ij}\rangle={1\over2}\int_Mv^{ij}v_{ij}\Phi^2dV(g).
\end{equation}
\begin{lemma}[gradient-flow formula for the lambda invariant]
  \label{lem:kl-ricci-flow-as-gradient-flow-gradient-flow-formula-lambda-invariant}
  \source{kleiner-lott}
  \dcref{Lemma 10.2}
  \group{Kleiner--Lott Entropy Functionals}
The formal gradient flow of $\lambda$ is
\begin{equation}
{\partial g_{ij}\over\partial t}=-2\left(R_{ij}-2\nabla_i\nabla_j\log\Phi\right).
\end{equation}
\end{lemma}
\begin{proof}
\sketch
\begin{equation}\label{eqnlambdadef}
\lambda(g)=\inf_{\int_Me^{-f}dV=1}{\mathcal F}(g,f).
\end{equation}
For the minimizing $f$, normalization gives
$\int_M(v/2-h)e^{-f}dV=0$. The eigenvalue equation yields
\begin{equation}\label{feqn}2\triangle f-|\nabla f|^2+R=\lambda.
\end{equation}
Thus (\ref{Fvariation}) reduces to
\begin{equation}\label{lambdavar}
\delta\lambda(v_{ij})=-\int_Me^{-f}v_{ij}(R_{ij}+\nabla_i\nabla_jf)dV,
\end{equation}
and
\begin{equation}\label{gradflow}
{\partial g_{ij}\over\partial t}=-2(R_{ij}+\nabla_i\nabla_jf)
=-2(R_{ij}-2\nabla_i\nabla_j\log\Phi).
\end{equation}
\end{proof}
\begin{equation}
\nabla_j\big((R_{ij}+\nabla_i\nabla_jf)e^{-f}\big)=0.
\end{equation}
Consequently the projection of (\ref{gradflow}) to ${\mathcal M}/\Diff(M)$ is
the projection of $g_t=-2\Ric$, so Ricci flow on the quotient is the gradient
flow of $\lambda$.

\subsection{The ${\mathcal W}$-functional}
\label{Basic Example for Section 3.1}

\begin{definition}[W-entropy definition]
  \label{def:kl-w-entropy-w-entropy-definition}
  \source{kleiner-lott}
  \dcref{Definition 11.1}
  \group{Kleiner--Lott Entropy Functionals}
\begin{equation}
{\mathcal W}(g,f,\tau)=\int_M\left[\tau(|\nabla f|^2+R)+f-n\right](4\pi\tau)^{-n/2}e^{-f}dV.
\end{equation}
\end{definition}

It is invariant under $\phi^*$ and under $(g,\tau)\mapsto(cg,c\tau)$. On flat
$\mathbb R^n$, $\tau=-t$, $f=|x|^2/(4\tau)$ satisfies (\ref{neweqn}), (\ref{Wnorm}),
and ${\mathcal W}=0$ by (\ref{Gaussianintegral})--(\ref{Gaussianintegral2}).

\subsection{Monotonicity of the ${\mathcal W}$-functional}
\dcref{I.3.1}

Put $\delta g_{ij}=v_{ij}$, $\delta f=h$, and $\delta\tau=\sigma$.
\begin{proposition}[first variation of W-entropy]
  \label{prop:kl-w-entropy-first-variation-w-entropy}
  \source{kleiner-lott}
  \dcref{Proposition 12.1}
  \group{Kleiner--Lott Entropy Functionals}
\begin{align}
\delta {\mathcal W}(v_{ij},h,\sigma)=\int_M\bigg[&\sigma(R+|\nabla f|^2)-\tau v_{ij}(R_{ij}+\nabla_i\nabla_jf)+h\\
&+\left\{\tau(2\triangle f-|\nabla f|^2+R)+f-n\right\}\left({v\over2}-h-{n\sigma\over2\tau}\right)\bigg](4\pi\tau)^{-n/2}e^{-f}dV.
\end{align}
\end{proposition}
\begin{proof}
\sketch
\begin{equation}
\delta((4\pi\tau)^{-n/2}e^{-f}dV)=\left({v\over2}-h-{n\sigma\over2\tau}\right)(4\pi\tau)^{-n/2}e^{-f}dV.
\end{equation}
Apply (\ref{Fvariation}) to the $\tau(R+|\nabla f|^2)$ term and use (\ref{id}).
\end{proof}

Fix $dm$ of mass $1$ with $(4\pi\tau)^{-n/2}e^{-f}dV=dm$. Then
\begin{equation}\label{Wvar}
\delta {\mathcal W}=\int_M[\sigma(R+|\nabla f|^2)-\tau v_{ij}(R_{ij}+\nabla_i\nabla_jf)+h](4\pi\tau)^{-n/2}e^{-f}dV.
\end{equation}
For
\begin{align}\label{oldeqn}
(g_{ij})_t&=-2(R_{ij}+\nabla_i\nabla_jf),\\
f_t&=-\triangle f-R+{n\over2\tau},\qquad\tau_t=-1,\notag
\end{align}
\begin{align}\label{Wtimeder}
{d{\mathcal W}\over dt}&=\int_M2\tau\left|R_{ij}+\nabla_i\nabla_jf-{1\over2\tau}g_{ij}\right|^2(4\pi\tau)^{-n/2}e^{-f}dV.
\end{align}
After a Lie derivative,
\begin{align}\label{neweqn}
(g_{ij})_t&=-2R_{ij},\\
f_t&=-\triangle f+|\nabla f|^2-R+{n\over2\tau},\qquad\tau_t=-1,\notag
\end{align}
and (\ref{Wtimeder}) remains valid, with
\begin{equation}\label{Wnorm}\int_M(4\pi\tau)^{-n/2}e^{-f}dV=1.\end{equation}
Define
\begin{equation}
\mu(g,\tau)=\inf\{{\mathcal W}(g,f,\tau):\int_M(4\pi\tau)^{-n/2}e^{-f}dV=1\}.
\end{equation}
For $\Phi=e^{-f/2}$ the variational problem is
\begin{equation}
(4\pi\tau)^{-n/2}\int_M[\tau(4|\nabla\Phi|^2+R\Phi^2)-2\Phi^2\log\Phi-n\Phi^2]dV,
\quad(4\pi\tau)^{-n/2}\int_M\Phi^2dV=1,
\end{equation}
with Euler--Lagrange equation
\begin{equation}
\tau(-4\triangle+R)\Phi=2\Phi\log\Phi+(\mu(g,\tau)+n)\Phi.
\end{equation}
The minimizer is positive and smooth, and $\mu(g(t),t_0-t)$ is nondecreasing.

\begin{definition}[shrinking breather definition]
  \label{def:kl-w-entropy-shrinking-breather-definition}
  \source{kleiner-lott}
  \dcref{Definition 12.17}
  \group{Kleiner--Lott Entropy Functionals}
A shrinking breather is a Ricci flow on $[t_1,t_2]$ with
$g(t_2)=c\phi^*g(t_1)$ for $c<1$ and $\phi\in\Diff(M)$.
\end{definition}
\begin{proposition}[rigidity of shrinking breathers]
  \label{prop:kl-w-entropy-rigidity-shrinking-breathers}
  \source{kleiner-lott}
  \dcref{Proposition 12.18}
  \group{Kleiner--Lott Entropy Functionals}
A shrinking breather is a gradient shrinking soliton.
\end{proposition}
\begin{proof}
\sketch
Set $t_0=(t_2-ct_1)/(1-c)$ and $\tau_i=t_0-t_i$, so $\tau_2=c\tau_1$. Scale and
diffeomorphism invariance give
\begin{equation}
\mu(g(t_2),\tau_2)=\mu((\tau_2/\tau_1)\phi^*g(t_1),\tau_2)=\mu(g(t_1),\tau_1),
\end{equation}
so equality holds in the monotonicity formula.
\end{proof}

\subsection{The no local collapsing theorem}
\dcref{I.4}
\begin{definition}[locally collapsing definition]
  \label{def:kl-no-local-collapsing-theorem-i-locally-collapsing-definition}
  \source{kleiner-lott}
  \dcref{Definition 13.1}
  \group{Kleiner--Lott Entropy Functionals}
A Ricci flow on $[0,T)$ is locally collapsing at $T$ if there are $t_k\to T$ and
$B_k=B(p_k,r_k)$ at time $t_k$ with $r_k^2/t_k$ bounded,
$|\Rm|(g(t_k))\le r_k^{-2}$ on $B_k$, and $r_k^{-n}\vol(B_k)\to0$.
\end{definition}
\begin{remark}[estimates for noncollapsing]
  \label{rem:kl-no-local-collapsing-theorem-i-estimates-noncollapsing}
  \source{kleiner-lott}
  \dcref{Remark 13.2}
  \group{Kleiner--Lott Entropy Functionals}
The ratio condition allows $T=\infty$; for finite $T$ it implies bounded $r_k$.
\end{remark}
\begin{theorem}[noncollapsing theorem]
  \label{thm:kl-no-local-collapsing-theorem-i-noncollapsing-theorem}
  \source{kleiner-lott}
  \dcref{Theorem 13.3}
  \group{Kleiner--Lott Entropy Functionals}
  \label{noncollthm} (cf. Theorem I.4.1)
If $M$ is closed and $T<\infty$, the flow is not locally collapsing at $T$.
\end{theorem}
\begin{proof}
\sketch
For a collapsing sequence use $\Phi_k=e^{-c_k/2}\phi(d_{t_k}(\cdot,p_k)/r_k)$, where
$\phi=1$ on $[0,1/2]$, $\phi=0$ on $[1,\infty)$, $0\le\phi\le1$, and
$|\phi'|\le10$. Normalize by
\begin{equation}\label{normalization}
e^{c_k}=\int_M(4\pi r_k^2)^{-n/2}\phi^2(d_{t_k}(x,p_k)/r_k)dV
\le(4\pi r_k^2)^{-n/2}\vol(B_k),
\end{equation}
so $c_k\to-\infty$. With $A_k(s)$ the distance-sphere mass and
\begin{equation}
\bar R_k(s)=r_k^2A_k(s)^{-1}\int_{S(p_k,r_ks)}R\,d\area,
\end{equation}
\begin{equation}\label{mess}
{\mathcal W}(g(t_k),f_k,r_k^2)=
{\int_0^1[4(\phi')^2+(\bar R_k+c_k-2\log\phi-n)\phi^2]A_k(s)ds\over\int_0^1\phi^2A_k(s)ds}.
\end{equation}
The gradient and logarithmic cutoff terms are bounded by
\begin{equation}
401\left({\int_0^1\sinh^{n-1}s\,ds\over\int_0^{1/2}\sinh^{n-1}s\,ds}-1\right),
\end{equation}
while $\bar R_k\le n(n-1)$; hence ${\mathcal W}\le\const+c_k\to-\infty$.
Monotonicity of $\mu$ contradicts continuity of $\mu(g(0),\tau)$.
\end{proof}
\begin{remark}[estimates for scalar curvature]
  \label{rem:kl-no-local-collapsing-theorem-i-estimates-scalar-curvature}
  \source{kleiner-lott}
  \dcref{Remark 13.13}
  \group{Kleiner--Lott Entropy Functionals}
The proof only requires $r_k^{-n}\vol(B(p_k,r_k))\to0$, an upper bound for
$r_k^2R$, and a uniform upper bound for
$\vol(B(p_k,r_k))/\vol(B(p_k,r_k/2))$. If the last ratio is at least $3^n$,
replace $r_k$ by $r_k/2$ and iterate; the normalized volume and scalar bound persist.
\end{remark}
\begin{definition}[kappa noncollapsing]
  \label{def:kl-no-local-collapsing-theorem-i-kappa-noncollapsing}
  \source{kleiner-lott}
  \dcref{Definition 13.14}
  \group{Kleiner--Lott Entropy Functionals}
  \label{noncollapsed1} (cf. Definition I.4.2)
$g$ is $\kappa$-noncollapsed on scale $\rho$ if every ball of radius $r<\rho$
with $|\Rm|\le r^{-2}$ throughout has volume at least $\kappa r^n$.
\end{definition}
\begin{remark}[structural remark on noncollapsing]
  \label{rem:kl-no-local-collapsing-theorem-i-structural-remark-noncollapsing}
  \source{kleiner-lott}
  \dcref{Remark 13.15}
  \group{Kleiner--Lott Entropy Functionals}
This is the slice version; the parabolic version used from I.7 onward is different.
\end{remark}

\subsection{The ${\mathcal W}$-functional as a time derivative} \label{I.5}
\dcref{I.5}
With $\tau=-t$ and the measure constraint in (\ref{oldeqn}),
\begin{align}
{d\over d\tau}\left(\tau\int_M(f-n/2)dm\right)&={\mathcal W}(g(t),f(t),\tau),
\end{align}
and under (\ref{neweqn}) the same identity holds with
$(4\pi\tau)^{-n/2}e^{-f}dV$. Likewise, under (\ref{newevolution}),
\begin{equation}
{d\over dt}\left(-\int_Mfe^{-f}dV\right)={\mathcal F}(g(t),f(t)).
\end{equation}

\section{Reduced length and reduced volume} \label{I.7}
\dcref{I.7}

Use backward time $\tau=t_0-t$, so $g_\tau=2\Ric$. For
$\gamma:[\tau_1,\tau_2]\to M$, $\tau_1\ge0$, define
\begin{equation}\label{defL}
{\mathcal L}(\gamma)=\int_{\tau_1}^{\tau_2}\sqrt\tau\left(R(\gamma(\tau))+|\dot\gamma(\tau)|^2\right)d\tau.
\end{equation}
If $X=\dot\gamma$, the Euler--Lagrange equation is
\begin{equation}
\nabla_XX-\tfrac12\nabla R+{1\over2\tau}X+2\Ric(X,\cdot)=0.
\end{equation}
For $p=\gamma(0)$, the initial vector is $v=\lim_{\tau\to0}\sqrt\tau X(\tau)$;
write ${\mathcal L}\exp_\tau(v)=\gamma(\tau)$. Let $L(q,\tau)$ be the infimal
${\mathcal L}$-length from $(p,0)$ to $(q,\tau)$ and set
\begin{equation}\label{reducedl}l(q,\tau)={L(q,\tau)\over2\sqrt\tau}.
\end{equation}
\begin{equation}\label{reducedv}\tilde V(\tau)=\int_M\tau^{-n/2}e^{-l(q,\tau)}dq.
\end{equation}

\subsection{Basic example for the reduced volume} \label{7Euclidean}
\dcref{I.7}
For static $\mathbb R^n$ with $p=0$,
\begin{equation}
\gamma(\tau)=\left({\tau\over\bar\tau}\right)^{1/2}q=2\sqrt\tau\,v,\qquad
L(q,\bar\tau)={|q|^2\over2\sqrt{\bar\tau}},\qquad l(q,\bar\tau)={|q|^2\over4\bar\tau},
\end{equation}
\begin{equation}
\bar L(q,\bar\tau)=2\sqrt{\bar\tau}L(q,\bar\tau)=|q|^2,qquad
\tilde V(\tau)=\int_{\mathbb R^n}\tau^{-n/2}e^{-|q|^2/(4\tau)}dq=(4\pi)^{n/2}.
\end{equation}

\subsection{Remarks about ${\mathcal L}$-geodesics and ${\mathcal L}\exp$} \label{Lremarks}
For a variation field $Y$ with $[X,Y]=0$,
\begin{equation}\label{firstvariation}
\delta_Y{\mathcal L}=2\sqrt\tau\langle X,Y\rangle\big|_{\tau_1}^{\tau_2}
+\int_{\tau_1}^{\tau_2}\sqrt\tau\langle Y,\nabla R-2\nabla_XX-4\Ric(X,\cdot)-X/\tau\rangle d\tau.
\end{equation}
Thus
\begin{equation}\label{Lgeod}\nabla_XX-\tfrac12\nabla R+{X\over2\tau}+2\Ric(X,\cdot)=0.\end{equation}
With $s=\sqrt\tau$ and $\hat X=d\gamma/ds=2sX$,
\begin{equation}\label{changetos}
{\mathcal L}=2\int_{s_1}^{s_2}\left({1\over4}|d\gamma/ds|^2+s^2R\right)ds.
\end{equation}
\begin{equation}\label{sE-L}
\nabla_{\hat X}\hat X-2s^2\nabla R+4s\Ric(\hat X,\cdot)=0.
\end{equation}
If curvature is bounded on compact time intervals and slices are complete, then
for every $v\in T_pM$ the corresponding ${\mathcal L}$-geodesic exists globally on
the interval, and ${\mathcal L}\exp_\tau:T_pM\to M$ is smooth. In particular,
\begin{equation}\label{dr}|\nabla R(x,\tau)|<{D\over\sqrt{\tau_2-\tau}}.\end{equation}
The minimizing ${\mathcal L}$-geodesics are smooth away from the ${\mathcal L}$-cut locus.
The cut locus has measure zero; on its complement ${\mathcal L}\exp_\tau$ is a
diffeomorphism from the minimizing initial-data domain $\Omega_\tau$, and
$\Omega_{\tau'}\subset\Omega_\tau$ for $\tau\le\tau'$.

\subsection{First derivatives of $L$}
\dcref{I.7.3--I.7.6}
Away from the cut locus, for the minimizing geodesic ending at $(q,\bar\tau)$,
\begin{equation}\label{grad1}\nabla L(q,\bar\tau)=2\sqrt{\bar\tau}X(\bar\tau),\end{equation}
\begin{equation}\label{gradsq}|\nabla L|^2=4\bar\tau|X|^2=-4\bar\tau R+4\bar\tau(R+|X|^2),\end{equation}
\begin{align}\label{timed}
L_{\bar\tau}&=2\sqrt{\bar\tau}R-\sqrt{\bar\tau}(R+|X|^2).
\end{align}
Set
\begin{equation}\label{Hexp}H(X)=-R_\tau-{R\over\tau}-2\langle\nabla R,X\rangle+2\Ric(X,X),qquad
K=\int_0^{\bar\tau}\tau^{3/2}H(X(\tau))d\tau.
\end{equation}
Then
\begin{equation}\label{grad2}
{dL(\gamma(\bar\tau),\bar\tau)\over d\bar\tau}
=\sqrt{\bar\tau}\left(R(\gamma(\bar\tau))+|X(\bar\tau)|^2\right).
\end{equation}
\begin{equation}\label{ddt}
{d\over d\tau}(R(\gamma(\tau))+|X(\tau)|^2)
=-H(X)-{1\over\tau}(R+|X|^2).
\end{equation}
\begin{equation}\label{parts}
\int_0^{\bar\tau}\tau^{3/2}{d\over d\tau}(R+|X|^2)d\tau=-K-L(q,\bar\tau).
\end{equation}
\begin{equation}\label{sub4}\bar\tau^{3/2}(R+|X|^2)=-K+\tfrac12L.\end{equation}
\begin{equation}\label{touse1}L_{\bar\tau}=2\sqrt{\bar\tau}R-{L\over2\bar\tau}+{K\over\bar\tau},qquad
\end{equation}
\begin{equation}\label{ttouse1}
|\nabla L|^2=-4\bar\tau R+{2L\over\sqrt{\bar\tau}}-{4K\over\sqrt{\bar\tau}}.
\end{equation}

\subsection{Second variation of ${\mathcal L}$}
\dcref{I.7.7}
For the evolving connection $g_\tau=2\Ric$,
\begin{equation}\label{xxx}
{d\over d\tau}\langle\nabla_YY,X\rangle
=\langle\nabla_X\nabla_YY,X\rangle+\langle\nabla_YY,\nabla_XX\rangle
+2\Ric(\nabla_YY,X)+2(\nabla_Y\Ric)(Y,X)-(\nabla_X\Ric)(Y,Y).
\end{equation}
For the test field satisfying (\ref{diffeq}),
\begin{equation}\label{sub1}
Q(\widetilde Y,\widetilde Y)=\int_0^{\bar\tau}\sqrt\tau\left[
\Hess_R(\widetilde Y,\widetilde Y)+2\langle R(\widetilde Y,X)\widetilde Y,X\rangle
+2\left|-\Ric(\widetilde Y,\cdot)+{\widetilde Y\over2\tau}\right|^2
-4(\nabla_{\widetilde Y}\Ric)(\widetilde Y,X)+2(\nabla_X\Ric)(\widetilde Y,\widetilde Y)\right]d\tau.
\end{equation}
\begin{equation}\label{sub2}
-2\sqrt{\bar\tau}\Ric(Y(\bar\tau),Y(\bar\tau))
=-2\int_0^{\bar\tau}{d\over d\tau}\left(\sqrt\tau\Ric(\widetilde Y,\widetilde Y)\right)d\tau.
\end{equation}
For $[X,Y]=0$, the quadratic form of the second variation is
\begin{align}\label{Hess}
Q(Y,Y)=\int_0^{\bar\tau}\sqrt\tau[\Hess_R(Y,Y)+2\langle R(Y,X)Y,X\rangle+2|\nabla_XY|^2\\
-4(\nabla_Y\Ric)(Y,X)+2(\nabla_X\Ric)(Y,Y)]d\tau.
\end{align}
Define
\begin{equation}
TY=-\nabla_X\nabla_XY-{1\over2\tau}\nabla_XY+\tfrac12\Hess_R(Y,\cdot)
-2(\nabla_Y\Ric)(X,\cdot)-2\Ric(\nabla_YX,\cdot).
\end{equation}
Then
\begin{equation}\label{Qeqn}Q(Y,Y)=2\int_0^{\bar\tau}\sqrt\tau\langle Y,TY\rangle d\tau
+2\sqrt{\bar\tau}\langle\nabla_XY(\bar\tau),Y(\bar\tau)\rangle.
\end{equation}
An ${\mathcal L}$-Jacobi field satisfies $TY=0$. If $q$ is outside the cut locus,
$Y(0)=0$, $Y(\bar\tau)=w$, then
\begin{equation}\label{hesseq}\Hess_L(w,w)=Q(Y,Y)=2\sqrt{\bar\tau}\langle\nabla_XY(\bar\tau),Y(\bar\tau)\rangle.
\end{equation}
For any field with the same endpoint values, $\Hess_L(w,w)\le Q(Y,Y)$.

\subsection{Hessian bound for $L$} \label{7.9}
\dcref{I.7.8--I.7.9}
For a unit endpoint vector, solve
\begin{equation}\label{diffeq}\nabla_X\widetilde Y=-\Ric(\widetilde Y,\cdot)+{1\over2\tau}\widetilde Y,qquad
\widetilde Y(\bar\tau)=Y(\bar\tau).
\end{equation}
Then $|\widetilde Y(\tau)|^2=\tau/\bar\tau$. The Hessian estimate is
\begin{equation}\label{sub3}
\Hess_L(Y(\bar\tau),Y(\bar\tau))\le {1\over\sqrt{\bar\tau}}-2\sqrt{\bar\tau}\Ric(Y,Y)
-\int_0^{\bar\tau}\sqrt\tau H(X,\widetilde Y)d\tau,
\end{equation}
where
\begin{align}\label{Htilde}
H(X,\widetilde Y)=&-\Hess_R(\widetilde Y,\widetilde Y)-2\langle R(\widetilde Y,X)\widetilde Y,X\rangle\\
&-4\left((\nabla_X\Ric)(\widetilde Y,\widetilde Y)-(\nabla_{\widetilde Y}\Ric)(\widetilde Y,X)\right)-2\Ric_\tau(\widetilde Y,\widetilde Y)\\
&+2|\Ric(\widetilde Y,\cdot)|^2-{1\over\tau}\Ric(\widetilde Y,\widetilde Y).
\end{align}

\subsection{The Laplacian of $L$} \label{7.10}
\dcref{I.7.10}
Summing (\ref{sub3}) over an orthonormal endpoint basis gives
\begin{equation}\label{touse2}
\triangle L\le {n\over\sqrt{\bar\tau}}-2\sqrt{\bar\tau}R-{K\over\bar\tau}.
\end{equation}

\subsection{Estimates on ${\mathcal L}$-Jacobi fields}
\dcref{I.7.11}
\begin{equation}\label{dY}
{d|Y|^2\over d\tau}\bigg|_{\bar\tau}\le {1\over\bar\tau}-{1\over\sqrt{\bar\tau}}\int_0^{\bar\tau}\sqrt\tau H(X,\widetilde Y)d\tau.
\end{equation}
Equality is equivalent to $\widetilde Y=Y$, in which case
\begin{equation}\label{sharp}{1\over\bar\tau}=2\Ric(Y,Y)+{1\over\sqrt{\bar\tau}}\Hess_L(Y,Y).
\end{equation}

\subsection{Monotonicity of the reduced volume $\widetilde V$} \label{monoV}
\begin{lemma}[estimates for reduced volume]
  \label{lem:kl-monotonicity-reduced-volume-estimates-reduced-volume}
  \source{kleiner-lott}
  \dcref{Lemma 23.1}
  \group{Kleiner--Lott Reduced Geometry}
  \label{lowerbound}
For $0<\bar\tau_1\le\bar\tau_2$, there are $C_1,C_2>0$ such that
\begin{equation}\label{lbound}l(q,\tau)\ge C_1d(p,q)^2-C_2\end{equation}
for $\tau\in[\bar\tau_1,\bar\tau_2]$.
\end{lemma}
\begin{proof}
\sketch
Using (\ref{changetos}), bounded curvature, and uniform metric comparability,
${\mathcal L}(\gamma)\ge\frac12\int|d\gamma/ds|^2ds-\const\ge\const d(p,q)^2-\const$.
\end{proof}

The change of variables on $\Omega_\tau$ gives
\begin{equation}
\tilde V(\tau)=\int_{T_pM}\tau^{-n/2}e^{-l({\mathcal L}\exp_\tau(v),\tau)}
{\mathcal J}(v,\tau)\chi_\tau(v)dv.
\end{equation}
For a fixed $v$, (\ref{sub4}) gives
\begin{equation}\label{bbound1}{dl(\gamma(\tau),\tau)\over d\tau}\bigg|_{\bar\tau}=-\tfrac12\bar\tau^{-3/2}K,
\end{equation}
and (\ref{dY}) gives
\begin{equation}\label{bbound2}{d\log{\mathcal J}(v,\tau)\over d\tau}\bigg|_{\bar\tau}
\le {n\over2\bar\tau}-\tfrac12\bar\tau^{-3/2}K.
\end{equation}
Thus the integrand and $\chi_\tau$ are nonincreasing, so $\tilde V$ is nonincreasing.
If equality holds, then
\begin{equation}2\Ric(\tau)+\tau^{-1/2}\Hess_{L(\tau)}=\tau^{-1}g(\tau),\end{equation}
and the flow is a gradient shrinking soliton.

\subsection{A differential inequality for $L$} \label{I.(7.15)}
\dcref{I.7.15}
With $\bar L=2\sqrt\tau L$, (\ref{touse1}) and (\ref{touse2}) imply
\begin{equation}\label{barrier}\bar L_{\bar\tau}+\triangle\bar L\le2n.
\end{equation}
This holds globally in the barrier sense: at a cut point, fix an initial segment
of a minimizing geodesic up to $\epsilon$ and minimize from its endpoint. The resulting
smooth upper barrier satisfies (\ref{barrier}) up to $o(1)$. The maximum principle
therefore makes $\min(\bar L-2n\bar\tau)$ nonincreasing.
\begin{lemma}[positivity for reduced l-geometry]
  \label{lem:kl-reduced-l-geometry-positivity-reduced-l-geometry}
  \source{kleiner-lott}
  \dcref{Lemma 24.3}
  \group{Kleiner--Lott Reduced Geometry}
  \label{neg}
For small positive $\bar\tau$, $\min\bar L(\cdot,\bar\tau)-2n\bar\tau<0$.
\end{lemma}
\begin{proof}
\sketch
The static curve at $p$ gives $\bar L=O(\bar\tau^2)$.
\end{proof}
Consequently $\min l(\cdot,\tau)\le n/2$. In the distributional sense,
\begin{equation}
l_{\bar\tau}-\triangle l+|\nabla l|^2-R+{n\over2\bar\tau}\ge0,
\end{equation}
equivalently
\begin{equation}\label{illustrate}(\partial_{\bar\tau}-\triangle)(\tau^{-n/2}e^{-l}dV)\le0.
\end{equation}
\begin{align}\label{pphieqn}
&\int_M\phi\,\bar\tau_2^{-n/2}e^{-l(\cdot,\bar\tau_2)}dV(\bar\tau_2)
-\int_M\phi\,\bar\tau_1^{-n/2}e^{-l(\cdot,\bar\tau_1)}dV(\bar\tau_1)\\
&\le\int_{\bar\tau_1}^{\bar\tau_2}\int_M(\triangle\phi)\,
\tau^{-n/2}e^{-l(\cdot,\tau)}dV(\tau)d\tau.
\end{align}
Using cutoffs and (\ref{lbound}) gives monotonicity of $\tilde V$. Also,
\begin{equation}\label{illustrate2}2\triangle l-|\nabla l|^2+R+{l-n\over\bar\tau}\le0.\end{equation}

\subsection{Estimates on the reduced length}
\dcref{I.7.2}
For nonnegative curvature operator and $\tau=\tau_0-t$,
\begin{equation}\label{ineqtoshow}H(X,Y)\ge-\Ric(Y,Y)\left({1\over\tau}+{1\over\tau_0-\tau}\right),\qquad
H(X)\ge-R\left({1\over\tau}+{1\over\tau_0-\tau}\right).
\end{equation}
For $\tau\le(1-c)\tau_0$,
\begin{equation}\label{I.(7.16)}|\nabla l|^2+R\le {C(c)l\over\tau},qquad
{d\over d\tau}\log|Y|^2\le {Cl+1\over\tau}.
\end{equation}

\subsection{The no local collapsing theorem II} \label{I.7.3}
\dcref{I.7.3}
\begin{definition}[kappa noncollapsing]
  \label{def:kl-no-local-collapsing-theorem-ii-kappa-noncollapsing}
  \source{kleiner-lott}
  \lean{KleinerLott.IsKappaNoncollapsedOnScale}
  \leanok
  \dcref{Definition 26.1}
  \group{Kleiner--Lott Reduced Geometry}
  \label{noncollapsed2}
A flow on $[0,T)$ is $\kappa$-noncollapsed on scale $\rho$ if, whenever
$r<\rho$, $t_0\ge r^2$, and $|\Rm|\le r^{-2}$ on
$B_{t_0}(x_0,r)\times[t_0-r^2,t_0]$, then
$\vol(B_{t_0}(x_0,r))\ge\kappa r^n$.
\end{definition}

\begin{theorem}[existence for kappa solutions]
  \label{thm:kl-no-local-collapsing-theorem-ii-existence-kappa-solutions}
  \source{kleiner-lott}
  \uses{def:kl-no-local-collapsing-theorem-ii-kappa-noncollapsing}
  \dcref{Theorem 26.2}
  \group{Kleiner--Lott Reduced Geometry}
  \label{nolocalcollapse}
For $n,T,\rho,K,c>0$, there is $\kappa=\kappa(n,K,c,\rho,T)>0$ such that any
complete flow on $[0,T)$ with compact-time curvature bounds,
$|\Rm|(g(0))\le K$, and $\inj(g(0))\ge c$ is $\kappa$-noncollapsed on scale $\rho$.
With the other parameters fixed, $\kappa$ is nonincreasing in $T$.
\end{theorem}
\begin{proof}
\sketch
If the theorem failed, choose collapsing balls with
$\epsilon_k=r_k^{-1}\vol(B_k)^{1/n}\to0$. The reduced volume at
$\epsilon_kr_k^2$ tends to zero: vectors with $|v|\le .1\epsilon_k^{-1/2}$ remain
inside the ball by
\begin{equation}\label{diffmess}
\left|{d\over d(\tau/(\epsilon_kr_k^2))}(\epsilon_k^{1/2}\sqrt\tau|X|)\right|
\le C\epsilon_k(\epsilon_k^{1/2}\sqrt\tau|X|)+C\epsilon_k^2\left({\tau\over\epsilon_kr_k^2}\right)^{1/2},
\end{equation}
while the remaining vectors are controlled by the Euclidean Gaussian bound.
The reduced volume at time $t_k$ has a uniform positive lower bound using
$\min l\le n/2$ and bounded geometry. This contradicts monotonicity.
\end{proof}

\section{Length distortion and forward noncollapsing}
\dcref{I.8}
\subsection{Length distortion estimates}
\label{secI.8.3}
\dcref{I.8.3}
\begin{lemma}[length distortion estimates lemma]
  \label{lem:kl-length-distortion-estimates-length-distortion-estimates-lemma}
  \source{kleiner-lott}
  \dcref{Lemma 27.1}
  \group{Kleiner--Lott Local Estimates}
  \label{crudelemma}
If $\Ric\le(n-1)K$, then
\begin{equation}{\dist_{t_1}(x_0,x_1)\over\dist_{t_0}(x_0,x_1)}\ge e^{-(n-1)K(t_1-t_0)}\quad(t_1>t_0).
\end{equation}
\end{lemma}
\begin{proof}
\sketch
For any curve, $dL/dt=-\int\Ric(\dot\gamma,\dot\gamma)/|\dot\gamma|\,ds\ge-(n-1)KL$;
integrate and minimize at $t_1$.
\end{proof}
\begin{remark}[structural remark on length distortion estimates]
  \label{rem:kl-length-distortion-estimates-structural-remark-length-distortion-estimates}
  \source{kleiner-lott}
  \dcref{Remark 27.5}
  \group{Kleiner--Lott Local Estimates}
  \label{similar}
If $\Ric\ge-(n-1)K$, then
\begin{equation}{\dist_{t_1}(x_0,x_1)\over\dist_{t_0}(x_0,x_1)}\le e^{(n-1)K(t_1-t_0)}.\end{equation}
\end{remark}
\begin{equation}\label{conclusion}{d\over dt}\dist_t(x_0,x_1)\ge-(n-1)K\dist_t(x_0,x_1).\end{equation}
\begin{lemma}[distance distortion lemma]
  \label{lem:kl-length-distortion-estimates-distance-distortion-lemma}
  \source{kleiner-lott}
  \dcref{Lemma 27.8}
  \group{Kleiner--Lott Local Estimates}
  \label{distancedist} (cf. Lemma I.8.3(b))
If $d_{t_0}(x_0,x_1)\ge2r_0$ and $\Ric(x,t_0)\le(n-1)K$ on the endpoint
$r_0$-balls, then
\begin{equation}{d\over dt}\dist_t(x_0,x_1)\big|_{t_0}\ge-2(n-1)\left({2\over3}Kr_0+r_0^{-1}\right).
\end{equation}
\end{lemma}
\begin{proof}
\sketch
Apply the second variation formula to the piecewise-linear endpoint cutoff field
$V_i=f e_i$, with $f=s/r_0$, $1$, and $(d-s)/r_0$ on the three subintervals;
$\int|\nabla_XV_i|^2=2/r_0$ and the endpoint curvature terms contribute at most
$2(n-1)K(2r_0/3)$.
\end{proof}
\begin{corollary}[estimates for length distortion estimates]
  \label{cor:kl-length-distortion-estimates-estimates-length-distortion-estimates}
  \source{kleiner-lott}
  \dcref{Corollary 27.16}
  \group{Kleiner--Lott Local Estimates}
  \label{bound1}
  \cite[Theorem 17.2]{Hamilton}
If $\Ric\le K$, $K>0$, then
\begin{equation}{d\over dt}\dist_t(x_0,x_1)\ge-\const(n)K^{1/2}.\end{equation}
\end{corollary}
\begin{proof}
  \uses{lem:kl-length-distortion-estimates-distance-distortion-lemma}
\sketch
Take $r_0=K^{-1/2}$ and use (\ref{conclusion}) for distances at most $2r_0$ and
Lemma \ref{distancedist} otherwise.
\end{proof}
\begin{lemma}[estimates for distance distortion]
  \label{lem:kl-length-distortion-estimates-estimates-distance-distortion}
  \source{kleiner-lott}
  \dcref{Lemma 27.18}
  \group{Kleiner--Lott Local Estimates}
  (cf. Lemma I.8.3(a)) \label{I.8.3a}
If $\Ric(x,t_0)\le(n-1)K$ on $B_{t_0}(x_0,r_0)$, then outside that ball
\begin{equation}d_t-\triangle d\ge-(n-1)\left({2\over3}Kr_0+r_0^{-1}\right)
\end{equation}
at $t_0$, in the barrier sense when necessary.
\end{lemma}

\subsection{No local collapsing propagates forward in time and to larger scales}
\dcref{I.8.2}
\begin{definition}[kappa solutions definition]
  \label{def:kl-no-local-collapsing-propagates-forward-time-larger-scales-kappa-solutions-definition}
  \source{kleiner-lott}
  \lean{KleinerLott.IsKappaCollapsedAt}
  \leanok
  \dcref{Definition 28.1}
  \group{Kleiner--Lott Local Estimates}
  (cf. Definition of I.8.1)
A flow is $\kappa$-collapsed at $(x_0,t_0)$ on scale $r$ if
$|\Rm|\le r^{-2}$ on $B_{t_0}(x_0,r)\times[t_0-r^2,t_0]$ but
$\vol(B_{t_0}(x_0,r))\le\kappa r^n$.
\end{definition}
\begin{theorem}[existence for kappa solutions]
  \label{thm:kl-no-local-collapsing-propagates-forward-time-larger-scales-existence-kappa-solutions}
  \source{kleiner-lott}
  \dcref{Theorem 28.2}
  \group{Kleiner--Lott Local Estimates}
  (cf. Theorem I.8.2) \label{8.2thm}
For each $A<\infty$ there is $\kappa(A)>0$ such that a complete flow on
$[0,r_0^2]$ with bounded sectional curvature, $\vol(B_0(x_0,r_0))\ge A^{-1}r_0^n$,
and $|\Rm|\le(nr_0^2)^{-1}$ on $B_0(x_0,r_0)\times[0,r_0^2]$ is not
$\kappa$-collapsed at time $r_0^2$ on scales below $r_0$ at points within
$Ar_0$ of $x_0$.
\end{theorem}
\begin{remark}[estimates for no local collapsing propagates forward in time and to larger scales]
  \label{rem:kl-no-local-collapsing-propagates-forward-time-larger-scales-estimates-no-local-collapsing-propagates-forward-time-larger-scales}
  \source{kleiner-lott}
  \dcref{Remark 28.3}
  \group{Kleiner--Lott Local Estimates}
  \label{8.2fix}
The factor $n$ in the curvature hypothesis provides bounded geometry in the
localized region used in the proof.
\end{remark}
\begin{proof}
  \uses{rem:kl-length-distortion-estimates-structural-remark-length-distortion-estimates,thm:kl-no-local-collapsing-theorem-ii-existence-kappa-solutions}
\sketch
Scale to $r_0=1$. Length distortion gives bounded geometry on the region
$\{\dist_{1/2}(y,x_0)\le1/10,\ 0\le t\le1/2\}$. Choose $\phi$ increasing with
\begin{equation}\label{phieqn}2(\phi')^2/\phi-\phi''\ge(2A+100n)\phi'-C(A)\phi.\end{equation}
For $h(y,t)=\phi(d(y,t)-A(2t-1))(\bar L(y,1-t)+2n+1)$, the barrier inequality gives
$\square h\ge-(2n+C(A))h$ at minima and hence
\begin{equation}h_{\min}(1/2)\le(2n+1)e^{n+C(A)/2}.
\end{equation}
This bounds $\min_{d(y,1/2)\le1/10}\bar L(y,1/2)$ and therefore gives a positive
lower bound for the reduced volume at time $1$.
\end{proof}

\section{Perelman's differential Harnack inequality}
\label{I.9}
\dcref{I.9}
For $\square=\partial_t-\triangle$, the formal adjoint is
$\square^*=-\partial_t-\triangle+R$. If
\begin{equation}u=(4\pi(T-t))^{-n/2}e^{-f},\qquad\square^*u=0.\end{equation}
set
\begin{equation}\label{veqn}v=[(T-t)(2\triangle f-|\nabla f|^2+R)+f-n]u.\end{equation}
\begin{proposition}[Harnack inequalities proposition]
  \label{prop:kl-perelman-s-differential-harnack-inequality-harnack-inequalities-proposition}
  \source{kleiner-lott}
  \dcref{Proposition 29.5}
  \group{Kleiner--Lott Differential Harnack}
  (cf. Proposition I.9.1) \label{localized}
\begin{equation}\label{I.(9.1)}
\square^*v=-2(T-t)\left|R_{ij}+\nabla_i\nabla_jf-{g_{ij}\over2(T-t)}\right|^2u.
\end{equation}
\end{proposition}
\begin{proof}
\sketch
\begin{equation}\label{der}{d\triangle\over dt}=2R_{ij}\nabla_i\nabla_j.
\end{equation}
The equation $\square^*u=0$ is equivalent to
$(\partial_t+\triangle)f=n/[2(T-t)]+|\nabla f|^2-R$. Substitution, the Ricci
flow evolution equations, and (\ref{der}) yield (\ref{I.(9.1)}).
\end{proof}
For compact $M$,
\begin{equation}
{d\over dt}{\mathcal W}(g,f,T-t)=2(T-t)\int_M\left|R_{ij}+\nabla_i\nabla_jf-{g_{ij}\over2(T-t)}\right|^2u,dV.
\end{equation}
\begin{corollary}[monotonicity for Harnack inequalities]
  \label{cor:kl-perelman-s-differential-harnack-inequality-monotonicity-harnack-inequalities}
  \source{kleiner-lott}
  \dcref{Corollary 29.19}
  \group{Kleiner--Lott Differential Harnack}
  (cf. Corollary I.9.2) \label{whenever}
If the maximum principle holds, $\max(v/u)$ is nondecreasing in $t$.
\end{corollary}
\begin{proof}
\sketch
\begin{equation}(\partial_t+\triangle)(v/u)=\frac{v\square^*u-u\square^*v}{u^2}-{2\over u}\langle\nabla u,\nabla(v/u)\rangle.
\end{equation}
Apply the maximum principle.
\end{proof}
\begin{corollary}[Harnack inequalities corollary]
  \label{cor:kl-perelman-s-differential-harnack-inequality-harnack-inequalities-corollary}
  \source{kleiner-lott}
  \dcref{Corollary 29.21}
  \group{Kleiner--Lott Differential Harnack}
  (cf. Corollary I.9.3) \label{corI.9.3}
If the flow is defined on $[0,T]$ and $u$ tends to a delta function at $T$, then
$v\le0$ for $t<T$.
\end{corollary}
\begin{proof}
\sketch
For positive $\square$-harmonic $h$, $d\int hv,dV/dt=-\int h\square^*v,dV\ge0$;
the terminal flat-space limit is zero, so positivity of $h$ implies $v\le0$.
\end{proof}
\begin{corollary}[convergence for Harnack inequalities]
  \label{cor:kl-perelman-s-differential-harnack-inequality-convergence-harnack-inequalities}
  \source{kleiner-lott}
  \dcref{Corollary 29.23}
  \group{Kleiner--Lott Differential Harnack}
  (cf. Corollary I.9.5) \label{compare}
If the terminal delta is at $(p,T)$, then $f(q,t)\le l(q,T-t)$.
\end{corollary}
\begin{proof}
  \uses{cor:kl-perelman-s-differential-harnack-inequality-monotonicity-harnack-inequalities}
\sketch
Equation (\ref{illustrate}) gives $\square^*((4\pi\tau)^{-n/2}e^{-l})\le0$.
Apply the previous maximum-principle argument to $e^{f-l}$ and use the terminal
flat-space limit.
\end{proof}
\begin{remark}[structural remark on Harnack inequalities]
  \label{rem:kl-perelman-s-differential-harnack-inequality-structural-remark-harnack-inequalities}
  \source{kleiner-lott}
  \uses{cor:kl-perelman-s-differential-harnack-inequality-convergence-harnack-inequalities}
  \dcref{Remark 29.24}
  \group{Kleiner--Lott Differential Harnack}
The same comparison follows from
\begin{equation}
{d\over d\tau}(\tau^{1/2}f(\gamma(\tau),\tau))\le\tfrac12\tau^{1/2}(R+|\dot\gamma|^2)
\end{equation}
by integrating from $(p,T)$.
\end{remark}

\section{Pseudolocality and future noncollapsing}
\dcref{I.10}
\subsection{The statement of the pseudolocality theorem}
\label{statement of I.10.1}
\begin{theorem}[pseudolocality theorem]
  \label{thm:kl-statement-pseudolocality-theorem-pseudolocality-theorem}
  \source{kleiner-lott}
  \dcref{Theorem 30.1}
  \group{Kleiner--Lott Pseudolocality}
  (cf. Theorem I.10.1) \label{thmI.10.1}
For every $\alpha>0$ there are $\delta,\epsilon>0$ such that a complete pointed
flow on $[0,(\epsilon r_0)^2]$ with
$R(x,0)\ge-r_0^{-2}$ and
\begin{equation}
\vol(\partial\Omega)^n\ge(1-\delta)c_n\vol(\Omega)^{n-1}
\end{equation}
for $x\in B_0(x_0,r_0)$ and $\Omega\subset B_0(x_0,r_0)$ satisfies
\begin{equation}
|\Rm|(x,t)<\alpha t^{-1}+(\epsilon r_0)^{-2}
\end{equation}
whenever $0<t\le(\epsilon r_0)^2$ and $d_t(x,x_0)\le\epsilon r_0$.
\end{theorem}

\subsection{Claim 1: A point selection argument}
\dcref{Claim 1 of I.10.1}
\begin{lemma}[a point selection argument lemma]
  \label{lem:kl-point-selection-argument-point-selection-argument-lemma}
  \source{kleiner-lott}
  \lean{KleinerLott.SmoothCompleteRicciFlowOn.exists_point_selection_of_smooth_complete_ricci_flow}
  \leanok
  \dcref{Claim 1}
  \group{Kleiner--Lott Pseudolocality}
  (cf. Claim 1 of I.10.1) \label{Claim 1}
If $g(\cdot)$ is a Ricci flow solution on $[0,\epsilon^2]$, $A>0$,
$A\epsilon<1/(100n)$, and
$|\Rm|(x,t)\ge\alpha t^{-1}+\epsilon^{-2}$ for some $(x,t)$ with
$0<t\le\epsilon^2$ and $d_t(x,x_0)\le\epsilon$, then there is
$(\bar x,\bar t)\in M_\alpha$ with $\bar t\in(0,\epsilon^2]$ and
$d_{\bar t}(\bar x,x_0)<(2A+1)\epsilon$ such that
\begin{equation}\label{I.(10.1)}|\Rm|(x',t')\le4|\Rm|(\bar x,\bar t)|\end{equation}
whenever
\begin{equation}\label{I.(10.2)}(x',t')\in M_\alpha,\quad t'\in(0,\bar t],\quad
d_{t'}(x',x_0)\le d_{\bar t}(\bar x,x_0)+A|\Rm|^{-1/2}(\bar x,\bar t).
\end{equation}
\end{lemma}
\begin{proof}
\sketch
Iterate a violating point, multiplying curvature by $4$ each step. Then
$|\Rm|(x_k,t_k)\ge4^{k-1}\epsilon^{-2}$ and
\begin{equation}d(x_k,t_k)\le\epsilon+2A|\Rm|^{-1/2}(x_1,t_1)\le(2A+1)\epsilon.
\end{equation}
Smoothness forces termination.
\end{proof}

\subsection{Claim 2: Getting parabolic regions}
\dcref{Claim 2 of I.10.1}
\begin{lemma}[getting parabolic regions lemma]
  \label{lem:kl-getting-parabolic-regions-getting-parabolic-regions-lemma}
  \source{kleiner-lott}
  \uses{lem:kl-point-selection-argument-point-selection-argument-lemma}
  \dcref{Claim 2}
  \group{Kleiner--Lott Pseudolocality}
  (cf. Claim 2 of I.10.1) \label{Claim 2}
For the point from Lemma \ref{Claim 1}, let $Q=|\Rm|(\bar x,\bar t)$. Then
\begin{equation}\label{I.(10.1a)}|\Rm|(x',t')\le4Q\end{equation}
whenever
\begin{equation}\label{I.(10.3)}\bar t-\tfrac12\alpha Q^{-1}\le t'\le\bar t,qquad
d_{\bar t}(x',\bar x)\le\tfrac1{10}AQ^{-1/2}.
\end{equation}
\end{lemma}
\begin{proof}
  \uses{lem:kl-length-distortion-estimates-distance-distortion-lemma,lem:kl-point-selection-argument-point-selection-argument-lemma}
\sketch
Points outside $M_\alpha$ satisfy $|\Rm|<\alpha/t'\le2Q$. For points in $M_\alpha$,
use (\ref{I.(10.2)}). Lemma \ref{distancedist} with
$r_0=AQ^{-1/2}/100$ gives
\begin{equation}\label{32.4}
d_{t'}(x_0,x')-d_{\bar t}(x_0,x')\le\tfrac12\alpha Q^{-1}2(n-1)\left({2\over3}4Q{AQ^{-1/2}\over100}+{100Q^{1/2}\over A}\right).
\end{equation}
For large $A$ this is at most $AQ^{-1/2}/2$, proving (\ref{I.(10.1a)}).
\end{proof}

\subsection{Claim 3: An upper bound on the integral of $v$}
\dcref{Claim 3 of I.10.1} \label{Claim3}
\begin{lemma}[integral bound in pseudolocality]
  \label{lem:kl-upper-bound-integral-integral-bound-pseudolocality}
  \source{kleiner-lott}
  \dcref{Lemma 33.1}
  \group{Kleiner--Lott Pseudolocality}
  \label{stoccomplete}
If pointed Ricci flows converge smoothly on compact subsets to a flow stochastically
complete for $\square^*$, their fundamental solutions converge smoothly on compact
subsets after passing to a subsequence.
\end{lemma}
\begin{proof}
\sketch
Use the uniform $L^1$ bound, parabolic regularity, and Fatou's lemma; stochastic
completeness forces the subsequential limit to equal the fundamental solution.
\end{proof}

\begin{lemma}[integral bound in pseudolocality]
  \label{lem:kl-upper-bound-integral-integral-bound-pseudolocality-2}
  \source{kleiner-lott}
  \dcref{Claim 3}
  \group{Kleiner--Lott Pseudolocality}
  (cf. Claim 3 of I.10.1) \label{Claim 3}
There is $\beta>0$ such that, for the fundamental solution based at
$(\bar x_k,\bar t_k)$ and $Q_k=|\Rm|(\bar x_k,\bar t_k)$, some
\begin{equation}
\tilde t_k\in[\bar t_k-\tfrac12\alpha Q_k^{-1},\bar t_k]
\end{equation}
satisfies
\begin{equation}\int_{B_k}v_kdV_k\le-\beta.\end{equation}
where $B_k$ is the time-$\tilde t_k$ ball centered at $\bar x_k$ of radius
$\sqrt{\bar t_k-\tilde t_k}$.
\end{lemma}
\begin{proof}
  \uses{cor:kl-perelman-s-differential-harnack-inequality-harnack-inequalities-corollary,lem:kl-getting-parabolic-regions-getting-parabolic-regions-lemma,lem:kl-upper-bound-integral-integral-bound-pseudolocality}
\sketch
Assume the contrary and rescale by $Q_k$. If injectivity radii stay positive,
Claim 2 gives a smooth limit with $|\Rm|\le4$ and $|\Rm|(\bar x_\infty,\bar t_\infty)=1$.
By (\ref{corI.9.3}), $v_\infty\le0$; vanishing integrals on the limit balls imply
$v_\infty=0$, hence
\begin{equation}\Ric(g_\infty)+\Hess f_\infty-{1\over2(\bar t_\infty-t)}g_\infty=0.
\end{equation}
The resulting curvature scaling forces flatness, contradiction. If injectivity
radii tend to zero, rescale to injectivity radius one and obtain a flat limit
$(-\infty,0]\times L$. Writing $L$ as a flat Euclidean bundle over a compact flat
base of dimension $n-m$, the fundamental solution satisfies
\begin{equation}v_\infty=(m-n)\left(1+\tfrac12\log(4\pi\tau)\right)u_\infty.\end{equation}
so the corresponding ball integrals tend to $-\infty$, again a contradiction.
\end{proof}

\subsection{Proof of the pseudolocality theorem}
\dcref{I.10.1} \label{pfI.10.1}
\begin{proof}
  \uses{cor:kl-perelman-s-differential-harnack-inequality-harnack-inequalities-corollary, lem:kl-getting-parabolic-regions-getting-parabolic-regions-lemma, lem:kl-upper-bound-integral-integral-bound-pseudolocality}
\sketch
Assume the theorem fails and choose a minimal counterexample sequence with
$\epsilon_k,\delta_k\to0$. Then
\begin{equation}\label{largebound}|\Rm|(x,t)<\alpha t_k^{-1}+2\epsilon_k^{-2}.\end{equation}
on $0<t\le\epsilon_k^2$, $d_t(x,x_0)\le\epsilon_k$. Apply Claims 1--3.
Choose $\phi=1$ on $(-\infty,1]$, $\phi=0$ on $[2,\infty)$, decreasing between,
with $\phi''\ge-10\phi'$ and $(\phi')^2\le10\phi$. Put
$\tilde d=d+200n\sqrt t$ and $h=\phi(\tilde d/(10A\epsilon))$. Lemma
\ref{I.8.3a} gives
\begin{equation}\square h\ge-(10A\epsilon)^{-2}10\phi,qquad
\int_Mhu\,dV\big|_{t=0}\ge1-A^{-2}.\label{ut}
\end{equation}
Using (\ref{localized}) and Claim 3,
\begin{equation}-\int_Mhv\,dV\big|_{t=0}\ge\beta(1-A^{-2}).\end{equation}
At $t=0$, set $\tilde u=hu$ and $\tilde f=f-\log h$. Integration by parts gives
\begin{equation}
\int(-2\triangle f+|\nabla f|^2)he^{-f}dV
=\int\left(-|\nabla\tilde f|^2+{|\nabla h|^2\over h^2}\right)he^{-f}dV,
\end{equation}
and hence, with $\bar t=\bar t_k$,
\begin{equation}\label{somme}
\beta(1-A^{-2})\le\int(-\bar t|\nabla\tilde f|^2-\tilde f+n)\tilde u,dV+(1+c)A^{-2}+\epsilon^2.
\end{equation}
After rescaling $\hat g=g/(2\bar t)$ and normalizing the resulting density $U_k$,
\begin{equation}\label{contra1}
\tfrac12\beta\le\int_{M_k}\left(-\tfrac12|\nabla F_k|^2-F_k+n\right)U_kd\hat V_k.
\end{equation}
Euclidean logarithmic Sobolev and symmetrization under the isoperimetric hypothesis
give
\begin{equation}\label{Beckner}\int_{\mathbb R^n}(-\tfrac12|\nabla F|^2-F+n)U\,dV\le0,
\end{equation}
and
\begin{equation}\label{contra2}
\int_{M_k}|\nabla F_k|^2U_kd\hat V_k\ge(1-\delta_k)^{2/n}n\exp\left(1-{2\over n}\int F_kU_kd\hat V_k\right).
\end{equation}
Together these imply
\begin{equation}
{n\over2}\left((1-\delta_k)^{2/n}e^x-1-x\right)\le-\tfrac12\beta,
\end{equation}
for $x=1-(2/n)\int F_kU_kd\hat V_k$, contradicting
$\lim_k\inf_{x\in\mathbb R}((1-\delta_k)^{2/n}e^x-1-x)=0$.
\end{proof}

\subsection{The volumes of future balls}
\dcref{I.10.2}
\begin{corollary}[volume corollary]
  \label{cor:kl-volumes-future-balls-volume-corollary}
  \source{kleiner-lott}
  \uses{thm:kl-statement-pseudolocality-theorem-pseudolocality-theorem}
  \dcref{Corollary 35.1}
  \group{Kleiner--Lott Pseudolocality}
  (cf. Corollary I.10.2) \label{corI.10.2}
Under the hypotheses of Theorem \ref{thmI.10.1},
\begin{equation}\vol_t(B_t(x,\sqrt t))\ge c(n)t^{n/2}\end{equation}
for $0<t\le(\epsilon r_0)^2$ and $x\in B_0(x_0,\epsilon r_0)$.
\end{corollary}

\subsection{$\kappa$-noncollapsing at future times}
\dcref{I.10.4}
\begin{corollary}[existence for kappa solutions]
  \label{cor:kl-noncollapsing-at-future-times-existence-kappa-solutions}
  \source{kleiner-lott}
  \dcref{Corollary 36.1}
  \group{Kleiner--Lott Pseudolocality}
  (cf. Corollary I.10.4)
There are $\delta,\epsilon>0$ such that for every $A>0$ there is $\kappa(A)>0$ with
the following property. Under the hypotheses of Theorem \ref{thmI.10.1}, if
$A^{-1}(\epsilon r_0)^2\le t\le(\epsilon r_0)^2$ and $d_t(x,x_0)\le Ar_0$,
then the flow is not $\kappa$-collapsed at $(x,t)$ on scales below $\sqrt t$.
\end{corollary}
\begin{proof}
  \uses{cor:kl-volumes-future-balls-volume-corollary,thm:kl-no-local-collapsing-propagates-forward-time-larger-scales-existence-kappa-solutions,thm:kl-statement-pseudolocality-theorem-pseudolocality-theorem}
Apply Theorem \ref{8.2thm} at time $A^{-1}(\epsilon r_0)^2$, using Theorem
\ref{thmI.10.1} and Corollary \ref{corI.10.2}.
\end{proof}

\subsection{Diffeomorphism finiteness}
\label{I.10.5}
\dcref{I.10.5}
\begin{theorem}[existence for Riemannian metrics]
  \label{thm:kl-diffeomorphism-finiteness-existence-riemannian-metrics}
  \source{kleiner-lott}
  \dcref{Theorem 37.1}
  \group{Kleiner--Lott Pseudolocality}
For each $n$ there is $\delta>0$ such that for all $r_0,V>0$ only finitely many
diffeomorphism types of compact $n$-manifolds satisfy
\begin{enumerate}
\item $R\ge-r_0^{-2}$;
\item $\vol(M,g_0)\le V$;
\item every domain in an $r_0$-ball satisfies
$\vol(\partial\Omega)^n\ge(1-\delta)c_n\vol(\Omega)^{n-1}$.
\end{enumerate}
\end{theorem}
\begin{proof}
  \uses{cor:kl-volumes-future-balls-volume-corollary,thm:kl-statement-pseudolocality-theorem-pseudolocality-theorem}
\sketch
Run Ricci flow. Pseudolocality gives existence through $(\epsilon r_0)^2$, a
two-sided sectional-curvature bound at that time, and Corollary \ref{corI.10.2}
gives a uniform lower bound for $\epsilon r_0$-balls. The scalar lower bound and
the maximum principle give a uniform upper volume bound at that time. A maximal
$2\epsilon r_0$-separated net therefore has uniformly bounded cardinality and
diameter. Standard finiteness for metrics with two-sided sectional-curvature
bounds, bounded diameter, and noncollapsed balls proves the theorem.
\end{proof}

\section{$\kappa$-solutions}
\label{secI.11.1}
\dcref{I.11.1--I.11.9}
\begin{definition}[kappa noncollapsing]
\label{def:kl-solutions-kappa-noncollapsing}
\source{kleiner-lott}
\dcref{Definition 38.1}
\group{Kleiner--Lott Kappa Solutions}

Given $\kappa>0$, a $\kappa$-solution is a Ricci flow on
 $(-\infty,C)$ or $(-\infty,C]$ with complete time slices, curvature
bounded on compact time intervals, nonnegative curvature operator, $R>0$,
and $\kappa$-noncollapsing at all scales.  The term may also refer to a
solution with this property for some $\kappa>0$.
\end{definition}
For ancient solutions $R_t\geq0$, so the two standard noncollapsing
definitions are equivalent up to changing $\kappa$ by distance distortion.
Finite-time blowup limits on compact manifolds are $\kappa$-solutions.

\subsection{Asymptotic solitons} \label{I.11.2}
\dcref{I.11.2}
Fix $(p,t_0)$ and define reduced length $l(q,\tau)$ and reduced volume
$\widetilde V(\tau)$ from $(p,t_0)$, with $\tau=t_0-t$; choose
$q(\tau)$ with $l(q(\tau),\tau)\leq n/2$.
\begin{proposition}[existence for Ricci solitons]
\label{prop:kl-asymptotic-solitons-existence-ricci-solitons}
\source{kleiner-lott}
\dcref{Proposition 39.1}
\group{Kleiner--Lott Kappa Solutions}

(cf. Proposition I.11.2) \label{asympsoliton}
There are $\overline\tau_i\to\infty$ such that the solutions on
$[t_0-\overline\tau_i,t_0-\overline\tau_i/2]$, rescaled by
$\overline\tau_i^{-1}$ at $(q(\overline\tau_i),t_0-\overline\tau_i)$,
converge on $[-1,-1/2]$ to a nonflat gradient shrinking soliton.
\end{proposition}
\begin{proof}
\uses{prop:kl-perelman-s-differential-harnack-inequality-harnack-inequalities-proposition}
\sketch The estimates
\begin{equation}
 |l^{1/2}(q,\tau)-l^{1/2}(q(\tau),\tau)|
 \leq\sqrt{C/(4\tau)}d_{t_0-\tau}(q,q(\tau)),\qquad \tau R\leq Cl
\end{equation}
and $\partial_\tau l\geq-(1+C)l/(2\tau)$ give
\begin{equation}\label{upperl}
l(q,\tau)\leq(\overline\tau/\tau)^{(1+C)/2}
\left(\sqrt{C/(4\overline\tau)}d_{t_0-\overline\tau}(q,q(\overline\tau))
+\sqrt{n/2}\right)^2.
\end{equation}
They give local curvature bounds after rescaling, and noncollapsing gives a
pointed limit.  The lower estimate
\begin{equation}\label{Perelmanbound}
l(q,\tau')\geq-l(q(\tau'))-1+C(n)d_{t_0-\tau'}(q,q(\tau'))^2/\tau'
\end{equation}
and polynomial volume growth give reduced-volume convergence.  Its limit is
a positive constant less than $(4\pi)^{n/2}$; equality in the reduced-length
inequality gives the soliton equation via Proposition \ref{localized}.
A flat limit would be $\mathbb R^n$ with reduced volume $(4\pi)^{n/2}$;
the lower estimate is the bound detailed in \cite{Ye1}.
\end{proof}

\subsection{Two dimensional $\kappa$-solutions} \label{I.11.3}
\dcref{I.11.3}
\begin{corollary}[kappa solutions corollary]
\label{cor:kl-two-dimensional-solutions-kappa-solutions-corollary}
\source{kleiner-lott}
\dcref{Corollary 40.1}
\group{Kleiner--Lott Kappa Solutions}

(cf. Corollary I.11.3) \label{2Dsoliton}
The only oriented two-dimensional $\kappa$-solution is the shrinking round
$2$-sphere.
\end{corollary}
\begin{proof}
\uses{prop:kl-asymptotic-solitons-existence-ricci-solitons}
\sketch The only nonflat oriented nonnegatively curved two-dimensional
gradient shrinking soliton is the round sphere; the noncompact case is
excluded using Proposition \ref{asympsoliton}, as in \cite{Ye2}.  The
asymptotic rescalings and Hamilton's surface pinching argument
\cite{Hamsurface} then make
every slice round.
\end{proof}
\begin{remark}[structural remark on Ricci solitons]
\label{rem:kl-two-dimensional-solutions-structural-remark-ricci-solitons}
\source{kleiner-lott}
\uses{cor:kl-two-dimensional-solutions-kappa-solutions-corollary}
\dcref{Remark 40.2}
\group{Kleiner--Lott Kappa Solutions}


The same classification can also be proved by the alternate argument in
Section \ref{altpf11.3}.
\end{remark}

\subsection{Asymptotic scalar curvature and asymptotic volume ratio}
\dcref{I.11.4}
\label{11.4proof}
\begin{definition}[asymptotic scalar curvature ratio definition]
\label{def:kl-asymptotic-scalar-curvature-asymptotic-volume-ratio-asymptotic-scalar-curvature-ratio-definition}
\source{kleiner-lott}
\dcref{Definition 41.1}
\group{Kleiner--Lott Kappa Solutions}

For a complete connected Riemannian manifold,
$\mathcal R=\limsup_{x\to\infty}R(x)d(x,p)^2$ is the asymptotic scalar
curvature ratio, independent of $p$.
\end{definition}
\begin{theorem}[compactness for kappa solutions]
\label{thm:kl-asymptotic-scalar-curvature-asymptotic-volume-ratio-compactness-kappa-solutions}
\source{kleiner-lott}
\dcref{Theorem 41.2}
\group{Kleiner--Lott Kappa Solutions}

\label{asympscalar}
If $(M,g(\cdot))$ is a noncompact $\kappa$-solution, then $\mathcal R=\infty$
on every time slice.
\end{theorem}
\begin{proof}
\sketch If $0<\mathcal R<\infty$, rescaling at points with
$R(x_k)d(x_k,p)^2\to\mathcal R$ gives an annular nonflat metric cone.  For
a radial bivector $P$, nonnegative curvature gives $\Rm P=0$; differentiating
twice and using
\begin{equation}\label{curvevol}
\Rm_t=\Delta\Rm+Q(\Rm)
\end{equation}
and $\Rm P=0$ gives
\begin{equation}\label{diff1}
(\nabla_{e_i}\Rm)P+\Rm(\nabla_{e_i}P)=0.
\end{equation}
Differentiating once more gives
\begin{equation}\label{diff2}
(\Delta\Rm)P+2\sum_i(\nabla_{e_i}\Rm)\nabla_{e_i}P
\Rm(\Delta P)=0.
\end{equation}
Thus $(\Delta\Rm)(P,P)>0$, so backward time has negative curvature.  If
$\mathcal R=0$, rescaling by $r_k^{-2}$, $r_k=d(p,x_k)$, gives a flat cone
away from the vertex.  Its unit sphere has principal curvatures $1$ and is
the standard sphere for $n\geq3$, making the original slice Euclidean
\cite[Theorem A]{Eschenburg}.  For $n=2$, the Harnack limit of the bounded
domains is a Euclidean unit disk, using
\cite[Theorem 10.8]{Burago-Gromov-Perelman},
whose boundary has length $2\pi$.  Both alternatives contradict $R>0$.
\end{proof}
\begin{lemma}[scalar curvature lemma]
\label{lem:kl-asymptotic-scalar-curvature-asymptotic-volume-ratio-scalar-curvature-lemma}
\source{kleiner-lott}
\dcref{Lemma 41.4}
\group{Kleiner--Lott Kappa Solutions}

The metric cone in the annular limit of the proof of
Theorem \ref{asympscalar} is smooth; its radial distance $\rho$ is smooth.
\end{lemma}
\begin{proof}
\sketch Along geodesics away from the vertex, $(\rho^2)''=2$ in the flat
two-dimensional cone containing the geodesic.  Rademacher's theorem and
exponential coordinates give $\operatorname{Hess}(\rho^2)=2g$ and hence
smoothness of $\rho$.
\end{proof}
\begin{definition}[asymptotic volume ratio definition]
\label{def:kl-asymptotic-scalar-curvature-asymptotic-volume-ratio-asymptotic-volume-ratio-definition}
\source{kleiner-lott}
\dcref{Definition 41.12}
\group{Kleiner--Lott Kappa Solutions}

For a complete $n$-manifold with nonnegative Ricci curvature,
$\mathcal V=\lim_{r\to\infty}r^{-n}\vol(B(p,r))$ is the asymptotic volume
ratio, independent of $p$.
\end{definition}
\begin{proposition}[existence for kappa solutions]
\label{prop:kl-asymptotic-scalar-curvature-asymptotic-volume-ratio-existence-kappa-solutions}
\source{kleiner-lott}
\dcref{Proposition 41.13}
\group{Kleiner--Lott Kappa Solutions}

(cf. Proposition I.11.4) \label{propI.11.4}
If $(M,g(\cdot))$ is noncompact, then $\mathcal V=0$ on every time slice.
There are $x_k\to\infty$ for which the flows based at $(x_k,t_0)$, rescaled
by $R(x_k,t_0)$, converge to a $\kappa$-solution splitting off $\mathbb R$.
\end{proposition}
\begin{proof}
\sketch If $\mathcal V>0$, choose $x_k,s_k$ with
$s_k/d(x_k,p)\to0$, $R(x_k)s_k^2\to\infty$, and $R\leq2R(x_k)$ on
$B(x_k,s_k)$, by the point-selection estimate \cite[Lemma 22.2]{Hamilton}.
The rescaled limit has positive asymptotic volume ratio and
a line, hence is $\mathbb R\times N$; dimension two is flat, and induction
rules out the lower-dimensional factor in higher dimensions.
\end{proof}

\subsection{Curvature and normalized volume control each other in a $\kappa$-solution}
\begin{corollary}[existence for kappa solutions]
\label{cor:kl-solution-curvature-normalized-volume-control-each-other-existence-kappa-solutions}
\source{kleiner-lott}
\dcref{Corollary 42.1}
\group{Kleiner--Lott Kappa Solutions}

\label{11.4cors}
\begin{enumerate}
\item $r_0^{-n}\vol(B(x_0,r_0))$ is bounded below iff $r_0^2R(x_0)$ is
bounded above.
\item Normalized volume is almost maximal iff $r_0^2R(x_0)$ is almost zero.
\item Pointed $\kappa$-solutions with a fixed-radius ball of controlled
normalized volume are precompact and limits are ancient, nonnegatively
curved, and $\kappa$-noncollapsed, with locally bounded slice curvature.
\item For $\eta=\eta(n,\kappa)$, $|\nabla R|\leq\eta R^{3/2}$ and
$|R_t|\leq\eta R^2$.  More generally,
\[
\left|\frac{\mathop{\rm D}^k}{\mathop{\rm D}t^k}\nabla^l\Rm\right|(y,t)
\leq C(n,\rho,k,l,\kappa)R(x,t)^{k+l/2+1}
\]
on $B_t(x,\rho R(x,t)^{-1/2})$.
\item There is $\alpha$ depending only on $\kappa$, with
$\alpha(s)\to\infty$ and
$R(y)d(x,y)^2\geq\alpha(R(x)d(x,y)^2)$.
\end{enumerate}
\end{corollary}
\begin{proof}
\uses{cor:kl-curvature-bounds-ricci-flow-solutions-nonnegative-curvature-operator-ass-existence-curvature-operators,cor:kl-length-distortion-estimates-estimates-length-distortion-estimates,cor:kl-necklike-behavior-at-infinity-three-dimensional-solution-strong-version-existence-kappa-solutions,cor:kl-necklike-behavior-at-infinity-three-dimensional-solution-weak-version-existence-kappa-solutions,cor:kl-two-dimensional-solutions-kappa-solutions-corollary,cor:kl-volume-bound-existence-curvature-operators,def:kl-notation-terminology-smooth-closeness-epsilon-necks,lem:kl-alternate-proof-corollary-2dsoliton-using-proposition-propi-11-4-corolla-existence-kappa-solutions,lem:kl-getting-parabolic-regions-getting-parabolic-regions-lemma,lem:kl-length-distortion-estimates-distance-distortion-lemma,lem:kl-necklike-behavior-at-infinity-three-dimensional-solution-strong-version-existence-kappa-solutions,lem:kl-point-selection-argument-point-selection-argument-lemma,prop:kl-asymptotic-scalar-curvature-asymptotic-volume-ratio-existence-kappa-solutions,thm:kl-compactness-space-three-dimensional-solutions-compactness-kappa-solutions,thm:kl-statement-pseudolocality-theorem-pseudolocality-theorem}
\sketch Point selection and relative volume comparison prove (1); the converse
uses a $c_n/2$-volume limit and Harnack.  The almost-maximal case is the
flat-limit version.  Rescaling to $R=1$ and applying precompactness proves
(3)--(5).
\end{proof}

\subsection{Alternate proof of Corollary \ref{2Dsoliton}}
\label{altpf11.3}
\begin{lemma}[existence for kappa solutions]
\label{lem:kl-alternate-proof-corollary-2dsoliton-using-proposition-propi-11-4-corolla-existence-kappa-solutions}
\source{kleiner-lott}
\dcref{Lemma 43.1}
\group{Kleiner--Lott Kappa Solutions}

\label{2dvol}
There is $v=v(\kappa)>0$ such that every two-dimensional $\kappa$-solution
satisfies
\begin{equation}\label{volnotsmall}
\vol(B_t(x,d(x,y)))\geq v d(x,y)^2.
\end{equation}
\end{lemma}
\begin{proof}
\uses{cor:kl-solution-curvature-normalized-volume-control-each-other-existence-kappa-solutions}
\sketch A violating sequence, midpoint rescaling, and
\begin{equation}\label{halfvol}
\bar r_k^{-2}\vol(B_{t_k}(z_k,\bar r_k))=\frac\pi2
\end{equation}
give a limit splitting off a
line but not equal to $\mathbb R^2$, hence a cylinder, contradicting
noncollapsing.
\end{proof}
\subsection{A volume bound}
\dcref{I.11.5}
\begin{corollary}[existence for curvature operators]
\label{cor:kl-volume-bound-existence-curvature-operators}
\source{kleiner-lott}
\dcref{Corollary 44.1}
\group{Kleiner--Lott Kappa Solutions}

(cf. Corollary I.11.5) \label{corI.11.5}
For every $\epsilon>0$ there is $A<\infty$ such that any sequence of
nonnegatively curved flows satisfying: $B(x_k,r_k)$ compactly contained at
time zero; $\frac12R\leq R(x_k,0)=Q_k$ on the ball and time interval;
$t_kQ_k\to-\infty$; and $r_k^2Q_k\to\infty$, obeys
\[
\vol(B(x_k,AQ_k^{-1/2}))\leq\epsilon(AQ_k^{-1/2})^n
\]
for all large $k$.
\end{corollary}
\begin{proof}
\uses{cor:kl-length-distortion-estimates-estimates-length-distortion-estimates,prop:kl-asymptotic-scalar-curvature-asymptotic-volume-ratio-existence-kappa-solutions}
\sketch Point selection and relative volume comparison give, for
$\widetilde r_k\leq\bar r_k$,
\begin{equation}\label{volbig}
\frac{\vol(B(y_k,\widetilde r_k))}{\widetilde r_k^n}
\geq\frac{\vol(B(x_k,r_k))}{(10r_k)^n}\geq\frac{c}{10^n}.
\end{equation}
Rescaling a contrary sequence gives an ancient nonflat limit with
positive asymptotic volume ratio, while Proposition \ref{propI.11.4} and
volume monotonicity force that ratio to vanish.
\end{proof}

\subsection{Curvature bounds for Ricci flow solutions with nonnegative curvature operator, assuming a lower volume bound}
\dcref{I.11.6}
\begin{corollary}[existence for curvature operators]
\label{cor:kl-curvature-bounds-ricci-flow-solutions-nonnegative-curvature-operator-ass-existence-curvature-operators}
\source{kleiner-lott}
\dcref{Corollary 45.1}
\group{Kleiner--Lott Kappa Solutions}

(cf. Corollary I.11.6) \label{corI.11.6}
For every $w>0$ there are $B,C<\infty$ and $\tau_0>0$ such that:
\begin{enumerate}
\item If $t_0\in[-r_0^2,0)$, the flow has nonnegative curvature operator on
$[t_0,0]$, $B_0(x_0,r_0)$ is compactly contained, and every time-$t$ ball
has volume at least $wr_0^n$, then
\begin{equation}\label{Rbound3}
R(x,t)\leq Cr_0^{-2}+B(t-t_0)^{-1}
\end{equation}
for $d_t(x,x_0)\leq r_0/4$.
\item On $[-\tau_0r_0^2,0]$, a time-zero volume bound
$\vol(B_0(x_0,r_0))\geq wr_0^n$ gives
$R(x,t)\leq Cr_0^{-2}+B(t+\tau_0r_0^2)^{-1}$ for $d_t(x,x_0)\leq r_0/4$.
\end{enumerate}
\end{corollary}
\begin{remark}[structural remark on curvature operators]
\label{rem:kl-curvature-bounds-ricci-flow-solutions-nonnegative-curvature-operator-ass-structural-remark-curvature-operators}
\source{kleiner-lott}
\uses{cor:kl-curvature-bounds-ricci-flow-solutions-nonnegative-curvature-operator-ass-existence-curvature-operators}
\dcref{Remark 45.4}
\group{Kleiner--Lott Kappa Solutions}


The proof of (a) uses a fixed bound $-t_0\leq cr_0^2$; here $c=1$.
This does not affect the later use of part (b).
\end{remark}
\begin{proof}
\proves{cor:kl-curvature-bounds-ricci-flow-solutions-nonnegative-curvature-operator-ass-existence-curvature-operators}
\uses{cor:kl-volume-bound-existence-curvature-operators,lem:kl-getting-parabolic-regions-getting-parabolic-regions-lemma,lem:kl-length-distortion-estimates-distance-distortion-lemma,lem:kl-point-selection-argument-point-selection-argument-lemma,thm:kl-statement-pseudolocality-theorem-pseudolocality-theorem}
\sketch Rescale to $r_0=1$, point-pick a first high-curvature point, and
integrate derivative bounds over its parabolic neighborhood.  Corollary
\ref{corI.11.5} contradicts unbounded curvature.  For (b), choose $\tau_0$
so volume distortion preserves a fixed fraction of the initial volume and
apply (a), using
\begin{equation}\label{cc}
R(x,t)\leq C(5^{-n}w)+B(5^{-n}w)(t+\tau)^{-1}.
\end{equation}
The corresponding volume comparison is
\begin{align}\label{lll}
5^{-n}w&\geq
\left(\frac14-10(n-1)(\tau\sqrt C+2\sqrt{B\tau})\right)^n w.
\end{align}
Choose $\tau_0$ so that
$\frac14-10(n-1)(\tau_0\sqrt C+2\sqrt{B\tau_0})=\frac15$.
\end{proof}
\begin{corollary}[existence for curvature operators]
\label{cor:kl-curvature-bounds-ricci-flow-solutions-nonnegative-curvature-operator-ass-existence-curvature-operators-3}
\source{kleiner-lott}
\dcref{Corollary 45.13}
\group{Kleiner--Lott Kappa Solutions}

(cf. end of Section I.11.6)\label{I.11.6end}
For every $w>0$ there are $B,C<\infty$ and $\tau_0>0$ such that, if the
curvature operator on each $B_t(x_0,r_0)$ for
$t\in[-\tau_0r_0^2,0]$ is at least $-r_0^{-2}$ and
$\vol(B_0(x_0,r_0))\geq wr_0^n$, then
\[
R(x,t)\leq Cr_0^{-2}+B(t+\tau_0r_0^2)^{-1}
\]
for $d_t(x,x_0)\leq r_0/4$.
\end{corollary}
\begin{proof}
\sketch The same blowup argument applies; the earlier volume comparison loses
only $e^{-\const\tau r_0^{-2}}$ under the lower curvature bound.
\end{proof}

\subsection{Compactness of the space of three-dimensional $\kappa$-solutions}
\dcref{I.11.7}
\label{pf11.7}
\begin{theorem}[compactness for kappa solutions]
\label{thm:kl-compactness-space-three-dimensional-solutions-compactness-kappa-solutions}
\source{kleiner-lott}
\dcref{Theorem 46.1}
\group{Kleiner--Lott Kappa Solutions}

(cf. Theorem I.11.7) \label{thmI.11.7}
For fixed $\kappa>0$, oriented three-dimensional $\kappa$-solutions are
compact modulo scaling.
\end{theorem}
\begin{proof}
\uses{cor:kl-solution-curvature-normalized-volume-control-each-other-existence-kappa-solutions,cor:kl-two-dimensional-solutions-kappa-solutions-corollary,prop:kl-asymptotic-scalar-curvature-asymptotic-volume-ratio-existence-kappa-solutions}
\sketch Normalize $R(x_k,0)=1$ and apply Corollary \ref{11.4cors}.  An
unbounded-curvature limit would contain necks with cross-sections tending to
zero, impossible by the Sharafutdinov exhaustion in the one-ended case and
by line splitting in the two-ended case.  The exhaustion is by totally convex
sets satisfying
\begin{equation}\label{C}
C_{t_1}=\{q\in C_{t_2}:\dist(q,\partial C_{t_2})\geq t_2-t_1\}.
\end{equation}
\end{proof}
\begin{remark}[compactness for kappa solutions]
\label{rem:kl-compactness-space-three-dimensional-solutions-compactness-kappa-solutions}
\source{kleiner-lott}
\dcref{Remark 46.3}
\group{Kleiner--Lott Kappa Solutions}

The compactness theorem applies to both compact and noncompact
$\kappa$-solutions.
\end{remark}
\begin{remark}[compactness for kappa solutions]
\label{rem:kl-compactness-space-three-dimensional-solutions-compactness-kappa-solutions-3}
\source{kleiner-lott}
\uses{cor:kl-two-dimensional-solutions-kappa-solutions-corollary}
\dcref{Remark 46.4}
\group{Kleiner--Lott Kappa Solutions}


The Ricci-flow hypothesis supplies compactness of the limiting
nonnegatively curved surface through the two-dimensional classification;
this is not automatic for a merely noncollapsed Riemannian manifold.
\end{remark}
\begin{corollary}[kappa solutions corollary]
\label{cor:kl-compactness-space-three-dimensional-solutions-kappa-solutions-corollary}
\source{kleiner-lott}
\dcref{Corollary 46.5}
\group{Kleiner--Lott Kappa Solutions}

Every asymptotic soliton of a three-dimensional $\kappa$-solution is itself
a $\kappa$-solution.
\end{corollary}

\subsection{Necklike behavior at infinity of a three-dimensional $\kappa$-solution - weak version}
\dcref{I.11.8}
\begin{definition}[kappa solution definition]
\label{def:kl-necklike-behavior-at-infinity-three-dimensional-solution-weak-version-kappa-solution-definition}
\source{kleiner-lott}
\uses{def:kl-notation-terminology-smooth-closeness-epsilon-necks}
\dcref{Definition 47.1}
\group{Kleiner--Lott Kappa Solutions}


Fix $\epsilon>0$.  With $Q=R(x_0,0)$, $x_0$ is an $\epsilon$-neck center if
\[
\{(x,t):-(\epsilon Q)^{-1}<t\leq0,\ d_0(x,x_0)^2<(\epsilon Q)^{-1}\}
\]
after scaling by $Q$ is $\epsilon$-close in the fixed smooth topology to
the corresponding evolving round-cylinder region.  Let $M_\epsilon$ be the
set of points that are not centers of $\epsilon$-necks.
\end{definition}
\begin{corollary}[existence for kappa solutions]
\label{cor:kl-necklike-behavior-at-infinity-three-dimensional-solution-weak-version-existence-kappa-solutions}
\source{kleiner-lott}
\dcref{Corollary 47.2}
\group{Kleiner--Lott Kappa Solutions}

(cf. Corollary I.11.8) \label{corI.11.8}
For every $\epsilon>0$ there is $C=C(\epsilon,\kappa)$ such that, for every
oriented noncompact three-dimensional $\kappa$-solution, $M_\epsilon$ is
compact, $\diam(M_\epsilon)\leq CQ^{-1/2}$, and
$C^{-1}Q\leq R(x,0)\leq CQ$ on $M_\epsilon$, where
$Q=R(x_0,0)$ for some $x_0\in\partial M_\epsilon$.
\end{corollary}
\begin{proof}
\uses{cor:kl-solution-curvature-normalized-volume-control-each-other-existence-kappa-solutions,prop:kl-asymptotic-scalar-curvature-asymptotic-volume-ratio-existence-kappa-solutions,thm:kl-compactness-space-three-dimensional-solutions-compactness-kappa-solutions}
\sketch A sequence in $M_\epsilon$ going to infinity gives, after scalar
curvature normalization, a line-splitting cylindrical limit, contradicting
the neck definition.  The same compactness argument at a point and a nearest
boundary point gives the diameter and curvature comparisons.
\end{proof}

\subsection{Necklike behavior at infinity of a three-dimensional $\kappa$-solution - strong version}
\dcref{I.11.8}
\label{strongversion}
\begin{corollary}[existence for kappa solutions]
\label{cor:kl-necklike-behavior-at-infinity-three-dimensional-solution-strong-version-existence-kappa-solutions}
\source{kleiner-lott}
\dcref{Corollary 48.1}
\group{Kleiner--Lott Kappa Solutions}

\label{neckstructure}
For every $\kappa>0$ there is $\epsilon_0>0$ such that, for
$0<\epsilon<\epsilon_0$, there is $\alpha(\epsilon,\kappa)$ and exactly one
of the following holds:
\begin{enumerate}
\item the flow is round cylindrical and every point is an $\epsilon$-neck;
\item $M$ is noncompact, $M_\epsilon\ne\varnothing$, and
$R(x)d(x,y)^2<\alpha$ for $x,y\in M_\epsilon$;
\item $M$ is compact, $x,y\in M_\epsilon$ satisfy $R(x)d(x,y)^2>\alpha$,
$M_\epsilon\subset B(x,\alpha R(x)^{-1/2})\cup
B(y,\alpha R(y)^{-1/2})$, and
$R(z)d(z,\overline{xy})^2<\alpha$ for $z\in M-M_\epsilon$;
\item $M$ is compact and some $x\in M_\epsilon$ satisfies
$R(x)d(x,z)^2<\alpha$ for all $z\in M$.
\end{enumerate}
\end{corollary}
\begin{proof}
\proves{cor:kl-necklike-behavior-at-infinity-three-dimensional-solution-strong-version-existence-kappa-solutions}
\sketch The all-neck case gives the cylinder by neck fibration and
noncollapsing.  If two non-neck points are $\alpha$-far, Lemma
\ref{annoy} gives alternative (3).  In the remaining compact case,
maximizing $R(x)d(x,\cdot)^2$ and examining a neck cross-section gives (4).
\end{proof}
\begin{lemma}[existence for kappa solutions]
\label{lem:kl-necklike-behavior-at-infinity-three-dimensional-solution-strong-version-existence-kappa-solutions}
\source{kleiner-lott}
\dcref{Lemma 48.3}
\group{Kleiner--Lott Kappa Solutions}

\label{annoy}
For every $\epsilon,\kappa>0$ there is $\alpha(\epsilon,\kappa)$ such that if
$x,y\in M_\epsilon$ and $R(x)d(x,y)^2>\alpha$, then for every $z$ either
$R(x)d(z,x)^2<\alpha$, or $R(y)d(z,y)^2<\alpha$, or
$R(z)d(z,\overline{xy})^2<\alpha$ and $z\notin M_\epsilon$.
\end{lemma}
\begin{proof}
\uses{cor:kl-solution-curvature-normalized-volume-control-each-other-existence-kappa-solutions,thm:kl-compactness-space-three-dimensional-solutions-compactness-kappa-solutions}
\sketch A violating sequence, rescaling, and nearest-point geometry produce
two rays with Tits angle at least $\pi/2$.  Cylindrical cross-sections force a
round cylindrical limit, contradicting the neck status; rescaling at the
nearest point makes $z$ an $\epsilon$-neck center.
\end{proof}
\subsection{More properties of $\kappa$-solutions}
\begin{proposition}[existence for kappa solutions]
\label{prop:kl-more-properties-solutions-existence-kappa-solutions}
\source{kleiner-lott}
\dcref{Proposition 49.1}
\group{Kleiner--Lott Kappa Solutions}

\label{neckcrit}
For all $\kappa,\alpha,\theta>0$ there is $\beta(\kappa,\alpha,\theta)<\infty$
such that $R(x)d(x,y_i)^2>\beta$ and
$\cangle_x(y_1,y_2)\geq\theta$ imply that $x$ is the center of an
$\alpha$-neck and $\cangle_x(y_1,y_2)\geq\pi-\alpha$.
\end{proposition}
\begin{proof}
\uses{lem:kl-necklike-behavior-at-infinity-three-dimensional-solution-strong-version-existence-kappa-solutions}
\sketch Compactness and line splitting give the neck; its cross-section
separates the far points and gives the angle estimate.
\end{proof}
\begin{corollary}[existence for kappa solutions]
\label{cor:kl-more-properties-solutions-existence-kappa-solutions}
\source{kleiner-lott}
\dcref{Corollary 49.2}
\group{Kleiner--Lott Kappa Solutions}

\label{longsegment}
For all $\kappa,\epsilon>0$ there is $\rho(\kappa,\epsilon)$ such that, if
$\eta$ is minimizing with endpoints $y_1,y_2$, $z'\in\eta$ is nearest to $z$,
and $R(z')d(z',y_i)^2>\rho$, then $z,z'$ are $\epsilon$-neck centers and
\[
\max\{R(z)d(z,z')^2,R(z')d(z,z')^2\}<4\pi^2.
\]
\end{corollary}
\begin{proof}
\sketch Apply the proposition at $z'$ with a smaller neck parameter; the
orthogonal segment $zz'$ lies near a spherical cross-section.
\end{proof}

\subsection{Getting a uniform value of $\kappa$}
\label{secI.11.9}
\dcref{I.11.9}
\begin{proposition}[existence for kappa solutions]
\label{prop:kl-getting-uniform-value-existence-kappa-solutions}
\source{kleiner-lott}
\dcref{Proposition 50.1}
\group{Kleiner--Lott Kappa Solutions}

\label{propI.11.9}
There is $\kappa_0>0$ such that every oriented three-dimensional
$\kappa$-solution is a $\kappa_0$-solution or a quotient of the round
shrinking $S^3$.
\end{proposition}
\begin{proof}
\uses{cor:kl-three-dimensional-noncompact-noncollapsed-gradient-shrinkers-standard-compactness-kappa-solutions,prop:kl-asymptotic-solitons-existence-ricci-solitons,thm:kl-no-local-collapsing-theorem-ii-existence-kappa-solutions}
\sketch Reduced-volume bounds for a collapsed unit-curvature ball are of
order $3(\kappa')^{3/2}$.  Asymptotic rescalings are compact round quotients
or the cylinder and its $\mathbb Z_2$ quotient; their volume lower bounds
exclude sufficiently small $\kappa'$, with the compact alternative as stated.
\end{proof}
\begin{remark}[existence for Ricci solitons]
\label{rem:kl-getting-uniform-value-existence-ricci-solitons}
\source{kleiner-lott}
\uses{cor:kl-three-dimensional-noncompact-noncollapsed-gradient-shrinkers-standard-compactness-kappa-solutions, prop:kl-asymptotic-solitons-existence-ricci-solitons, thm:kl-compactness-space-three-dimensional-solutions-compactness-kappa-solutions}
\dcref{Remark 50.2}
\group{Kleiner--Lott Kappa Solutions}


In dimension three, Theorem \ref{thmI.11.7} supplies the global curvature
bound required when applying the noncompact shrinker classification to an
asymptotic soliton.
\end{remark}

\subsection{Three-dimensional noncompact $\kappa$-noncollapsed gradient shrinkers are standard}
\dcref{II.1.2}
\label{II.1.2}
\begin{lemma}[existence for kappa solutions]
\label{lem:kl-three-dimensional-noncompact-noncollapsed-gradient-shrinkers-standard-existence-kappa-solutions}
\source{kleiner-lott}
\dcref{Lemma 51.1}
\group{Kleiner--Lott Kappa Solutions}

\label{nosoliton} (cf. Lemma of II.1.2)
There is no complete oriented noncompact three-dimensional
$\kappa$-noncollapsed gradient shrinking soliton with bounded positive
sectional curvature.
\end{lemma}
\begin{proof}
\sketch The shrinker equation and its derivative are
\begin{equation}\label{IIeqn}
\nabla_i\nabla_jf+R_{ij}+\frac1{2t}g_{ij}=0,\qquad
\nabla_iR=2R_{ij}\nabla_jf.\label{gradR}
\end{equation}
Ricci integral estimates make $\nabla f$ asymptotically radial and its flow
strictly increases $R$.  Splitting limits give $R\to1$ at time $-1$.
For far level surfaces $N$,
\[
R^N=R-2\Ric(X,X)+2\det S,\qquad
\det(\Hess f|_{TN})\leq\tfrac14(1-R+\Ric(X,X))^2.
\]
Thus $R^N<1$ and the levels are convex with nondecreasing area.  Cylindrical
limits give area at most $8\pi$, while Gauss--Bonnet gives
$\int_NR^N\,dA=8\pi$, a contradiction.
\end{proof}
\begin{corollary}[compactness for kappa solutions]
\label{cor:kl-three-dimensional-noncompact-noncollapsed-gradient-shrinkers-standard-compactness-kappa-solutions}
\source{kleiner-lott}
\dcref{Corollary 51.22}
\group{Kleiner--Lott Kappa Solutions}

\label{cornosoliton}
The only complete oriented noncompact three-dimensional $\kappa$-noncollapsed
gradient shrinking solitons with bounded nonnegative sectional curvature are
the evolving round $\mathbb R\times S^2$ and its
$\mathbb Z_2$ quotient $\mathbb R\times_{\mathbb Z_2}S^2$.
\end{corollary}
\begin{proof}
\uses{cor:kl-two-dimensional-solutions-kappa-solutions-corollary,lem:kl-three-dimensional-noncompact-noncollapsed-gradient-shrinkers-standard-existence-kappa-solutions,thm:kl-maximum-principles-existence-curvature-operators}
\sketch The lemma rules out positive sectional curvature.  Splitting gives a
line in the universal cover, the two-dimensional classification gives the
standard cylinder, and noncollapsing leaves only it and its $\mathbb Z_2$
quotient.
\end{proof}
\section{Canonical neighborhoods and curvature bounds}
\dcref{I.12}
\subsection{Canonical neighborhood theorem}
\label{I.12.1}
\dcref{I.12.1}
\begin{definition}[almost nonnegative curvature]
\label{def:kl-canonical-neighborhood-theorem-almost-nonnegative-curvature}
\source{kleiner-lott}
\dcref{Definition 52.1}
\group{Kleiner--Lott Canonical Neighborhoods}

(cf. I.12) \label{pinchingdef}
Let $\Phi\in C^\infty(\mathbb R)$ be positive and nondecreasing, with
$\Phi(s)/s$ decreasing to zero for $s>0$.  A Ricci flow has
$\Phi$-almost nonnegative curvature if
\begin{equation}\label{phieqnn}
\Rm(x,t)\geq-\Phi(R(x,t)).
\end{equation}
\end{definition}
\begin{remark}[existence for canonical neighborhoods]
\label{rem:kl-canonical-neighborhood-theorem-existence-canonical-neighborhoods}
\source{kleiner-lott}
\uses{def:kl-canonical-neighborhood-theorem-almost-nonnegative-curvature}
\dcref{Remark 52.3}
\group{Kleiner--Lott Canonical Neighborhoods}


The condition implies $R\geq-6\Phi(0)$ and, since sectional curvatures sum
to $R$,
\begin{equation}\label{double}
-\Phi(R)\leq\Rm\leq\frac R2+
\left(\frac{n(n-1)}2-1\right)\Phi(R).
\end{equation}
It is the formulation supplied by the three-dimensional Hamilton--Ivey
pinching estimate.
\end{remark}
\begin{lemma}[convergence for canonical neighborhoods]
\label{lem:kl-canonical-neighborhood-theorem-convergence-canonical-neighborhoods}
\source{kleiner-lott}
\dcref{Lemma 52.5}
\group{Kleiner--Lott Canonical Neighborhoods}

\label{almost nonneg}
If complete pointed manifolds with $\Phi$-almost nonnegative curvature,
rescaled by $Q_k\to\infty$, converge smoothly, then the limit has
nonnegative sectional curvature.
\end{lemma}
\begin{proof}
\uses{def:kl-canonical-neighborhood-theorem-almost-nonnegative-curvature}
\sketch If the approximating scalar curvatures stay bounded, \eqref{double}
makes the rescaled curvature vanish.  If they diverge, write
$Q_k^{-1}\Rm=(Q_k^{-1}R)(\Rm/R)$ and use
$\Rm/R\geq-\Phi(R)/R\to0$.
\end{proof}
\begin{theorem}[compactness for canonical neighborhoods]
\label{thm:kl-canonical-neighborhood-theorem-compactness-canonical-neighborhoods}
\source{kleiner-lott}
\dcref{Theorem 52.7}
\group{Kleiner--Lott Canonical Neighborhoods}

(cf. Theorem I.12.1) \label{thmI.12.1}
Given $\epsilon,\kappa,\sigma>0$ and $\Phi$ as above, there is $r_0>0$ such
that a complete three-dimensional Ricci flow on $[0,T]$, $T\geq1$, with
bounded curvature on compact time intervals, $\Phi$-almost nonnegative
curvature, and $\kappa$-noncollapsing below scale $\sigma$ has this
property: if $t_0\geq1$ and $Q=R(x_0,t_0)\geq r_0^{-2}$, then
\[
\{(x,t):d_{t_0}(x,x_0)^2<(\epsilon Q)^{-1},\
t_0-(\epsilon Q)^{-1}\leq t\leq t_0\},
\]
after scaling by $Q$, is $\epsilon$-close to the corresponding region of a
$\kappa$-solution.
\end{theorem}
\begin{proof}
\sketch Suppose the theorem fails and rescale a sequence of bad points to
unit scalar curvature.  Point-picking produces points for which all nearby
points of substantially larger curvature are good.  The rescaled time slices
have bounded curvature on each fixed ball by the maximal-radius argument:
otherwise a further limit ends on a nonflat metric cone, contradicting
Theorem \ref{asympscalar}.  Lemmas \ref{rescaled} and
\ref{almost nonneg} then give a backward nonnegatively curved limit, which
extends to an ancient $\kappa$-noncollapsed flow.  It is a
$\kappa$-solution, contradicting the selected points' badness.
\begin{equation}\label{addbd}
|d_t(x,y)-d_{t_0}(x,y)|<C
\end{equation}
on the limiting backward interval.  In the noncompact case, asymptotic
conical geometry supplies $D$ such that, for $d_{t_0}(y,x_0)>D$, some $x$
satisfies
\begin{equation}\label{dists}
d_{t_0}(x,y)=d_{t_0}(y,x_0),\qquad
d_{t_0}(x,x_0)\geq\frac32d_{t_0}(y,x_0).
\end{equation}
\end{proof}
\begin{remark}[existence for canonical neighborhoods]
\label{rem:kl-canonical-neighborhood-theorem-existence-canonical-neighborhoods-5}
\source{kleiner-lott}
\uses{rem:kl-canonical-neighborhood-theorem-structural-remark-canonical-neighborhoods, thm:kl-canonical-neighborhood-theorem-compactness-canonical-neighborhoods}
\dcref{Remark 52.8}
\group{Kleiner--Lott Canonical Neighborhoods}


\label{wangg}
The theorem permits noncompact flows and assumes noncollapsing below a fixed
scale $\sigma$; the source formulation uses a scale tied to $r_0$.
\end{remark}
\begin{lemma}[canonical neighborhoods lemma]
\label{lem:kl-canonical-neighborhood-theorem-canonical-neighborhoods-lemma}
\source{kleiner-lott}
\uses{thm:kl-canonical-neighborhood-theorem-compactness-canonical-neighborhoods}
\dcref{Lemma 52.9}
\group{Kleiner--Lott Canonical Neighborhoods}


\label{ptpicking}
There are $H_k\to\infty$, $x_k$, $t_k\geq1/2$, and
$Q_k=R(x_k,t_k)\to\infty$ such that Theorem \ref{thmI.12.1} fails at
$(x_k,t_k)$ but holds at every $(y,t)$ in
$M_k\times[t_k-H_kQ_k^{-1},t_k]$ with $R(y,t)\geq2R(x_k,t_k)$.
\end{lemma}
\begin{proof}
\uses{thm:kl-canonical-neighborhood-theorem-compactness-canonical-neighborhoods}
\sketch Choose $H_k$ with
$H_kR(\widehat x_k,\widehat t_k)^{-1}\leq1/10$ and repeatedly replace a bad
point by a bad point of at least twice its scalar curvature in the indicated
backward interval; termination gives the claim.
\end{proof}
\begin{lemma}[existence for canonical neighborhoods]
\label{lem:kl-canonical-neighborhood-theorem-existence-canonical-neighborhoods}
\source{kleiner-lott}
\dcref{Lemma 52.10}
\group{Kleiner--Lott Canonical Neighborhoods}

\label{grad}
If $R(x,t)>0$ and the parabolic neighborhood
$B_t(x,(\epsilon R(x,t))^{-1/2})\times
[t-(\epsilon R(x,t))^{-1},t]$ is $\epsilon$-close to a
$\kappa$-solution region, then for $C=C(\kappa)$,
\[
|\nabla R^{-1/2}|(x,t)\leq C,\qquad|\partial_tR^{-1}|(x,t)\leq C.
\]
\end{lemma}
\begin{proof}
\uses{thm:kl-compactness-space-three-dimensional-solutions-compactness-kappa-solutions}
\sketch Rescale to unit scalar curvature and apply Theorem \ref{thmI.11.7}.
\end{proof}
\begin{lemma}[canonical neighborhoods lemma]
\label{lem:kl-canonical-neighborhood-theorem-canonical-neighborhoods-lemma-8}
\source{kleiner-lott}
\dcref{Lemma 52.11}
\group{Kleiner--Lott Canonical Neighborhoods}

(cf. Claim 1 of I.12.1) \label{claim1}
For $t_k-\frac12H_kQ_k^{-1}\leq\overline t\leq t_k$, put
$\overline Q_k=Q_k+|R_k(\overline x,\overline t)|$.  For some
$c=c(\kappa)>0$,
\[
R_k(x,t)\leq4\overline Q_k
\]
whenever $\overline t-c\overline Q_k^{-1}\leq t\leq\overline t$ and
$d_{\overline t}(x,\overline x)\leq c\overline Q_k^{-1/2}$.
\end{lemma}
\begin{proof}
\uses{lem:kl-canonical-neighborhood-theorem-canonical-neighborhoods-lemma,lem:kl-canonical-neighborhood-theorem-existence-canonical-neighborhoods}
\sketch Join $(x,t)$ to $(\overline x,\overline t)$ by a time segment and a
minimizing geodesic; stop at curvature $2Q_k$ and integrate Lemma
\ref{grad}.
\end{proof}
\begin{lemma}[canonical neighborhoods lemma]
\label{lem:kl-canonical-neighborhood-theorem-canonical-neighborhoods-lemma-9}
\source{kleiner-lott}
\dcref{Lemma 52.12}
\group{Kleiner--Lott Canonical Neighborhoods}

\label{rescaled}
For the $Q_k$-rescaled flows, put
$\widetilde Q_k=1+|\overline R_k(\overline x,\overline t)|$.  Then
$\overline R_k(x,t)\leq4\widetilde Q_k$ if
$\overline t-c\widetilde Q_k^{-1}\leq t\leq\overline t$ and
$d_{\overline t}(x,\overline x)\leq c\widetilde Q_k^{-1/2}$.
\end{lemma}
\begin{proof}
\uses{lem:kl-canonical-neighborhood-theorem-canonical-neighborhoods-lemma-8}
\sketch This is Lemma \ref{claim1} after parabolic rescaling.
\end{proof}
\begin{lemma}[existence for canonical neighborhoods]
\label{lem:kl-canonical-neighborhood-theorem-existence-canonical-neighborhoods-10}
\source{kleiner-lott}
\dcref{Lemma 52.14}
\group{Kleiner--Lott Canonical Neighborhoods}

\label{nearlyround}
In the maximal-radius argument, let $Z$ be the limit, $y_\infty$ the
curvature blowup point, and $\gamma_\infty$ the limiting segment.  There are
functions $c$ and $\epsilon'$ depending only on $\kappa$, with
$c(s)\to\infty$ and $\epsilon'(s)\to0$ as $s\to0$, such that:
\begin{enumerate}
\item $\overline R_\infty(w)d(y_\infty,w)^2>c(\epsilon)$ for
$w\in\gamma_\infty$;
\item if $w$ is sufficiently close to $y_\infty$, then
$(Z,w,\overline R_\infty(w)\overline g_\infty)$ is
$2\epsilon'(\epsilon)$-close in $C^2$ to a round cylinder.
\end{enumerate}
\end{lemma}
\begin{proof}
\uses{lem:kl-canonical-neighborhood-theorem-canonical-neighborhoods-lemma,lem:kl-canonical-neighborhood-theorem-existence-canonical-neighborhoods,prop:kl-more-properties-solutions-existence-kappa-solutions}
\sketch Point-picking supplies close $\kappa$-solution models along the
segment; Lemma \ref{grad} makes curvature diverge at the terminal point and
Proposition \ref{neckcrit} makes the models cylindrical.  Their cylindrical
balls form a nonnegatively curved Alexandrov limit whose tangent cone gives
the annular-cone contradiction of Theorem \ref{asympscalar}.
\end{proof}
\begin{remark}[structural remark on canonical neighborhoods]
\label{rem:kl-canonical-neighborhood-theorem-structural-remark-canonical-neighborhoods}
\source{kleiner-lott}
\uses{rem:kl-canonical-neighborhood-theorem-existence-canonical-neighborhoods-5, thm:kl-canonical-neighborhood-theorem-compactness-canonical-neighborhoods, thm:kl-no-local-collapsing-theorem-ii-existence-kappa-solutions}
\dcref{Remark 52.19}
\group{Kleiner--Lott Canonical Neighborhoods}


\label{wang}
The fixed-scale noncollapsing hypothesis is what ensures that the ancient
limit constructed in Theorem \ref{thmI.12.1} is a $\kappa$-solution; the
source formulation alone yields only a small-scale noncollapsing conclusion.
\end{remark}

\subsection{Later scalar curvature bounds on bigger balls from curvature and volume bounds}
\dcref{I.12.2}
\label{I.12.2}
\begin{theorem}[positivity for scalar curvature]
\label{thm:kl-later-scalar-curvature-bounds-bigger-balls-from-curvature-volume-bounds-positivity-scalar-curvature}
\source{kleiner-lott}
\dcref{Theorem 53.1}
\group{Kleiner--Lott Canonical Neighborhoods}

\label{thmI.12.2}
For every $A<\infty$ there are $K(A)<\infty$ and $\rho(A)>0$ such that a
three-dimensional $\Phi$-almost nonnegatively curved flow defined on
$[0,r_0^2]$, satisfying $r_0^2\Phi(r_0^{-2})<\rho$,
$|\Rm|\leq(3r_0^2)^{-1}$ on the time-zero $r_0$-ball throughout the interval,
and $\vol(B_0(x_0,r_0))\geq A^{-1}r_0^3$, satisfies
\[
R(x,r_0^2)\leq K r_0^{-2}\qquad\text{when }d_{r_0^2}(x,x_0)<Ar_0.
\]
\end{theorem}
\begin{proof}
\uses{lem:kl-canonical-neighborhood-theorem-canonical-neighborhoods-lemma-9,lem:kl-length-distortion-estimates-distance-distortion-lemma,thm:kl-canonical-neighborhood-theorem-compactness-canonical-neighborhoods,thm:kl-no-local-collapsing-propagates-forward-time-larger-scales-existence-kappa-solutions}
\sketch Lemma \ref{12.2lemma} models every unbounded final-time region on a
$\kappa$-solution.  The maximal rescaled-ball limit is nonnegatively curved
by the pinching estimate and Lemma \ref{almost nonneg}; compactness and the
ancient extension contradict the assumed failure of the estimate.
\end{proof}
\begin{remark}[existence for noncollapsing]
\label{rem:kl-later-scalar-curvature-bounds-bigger-balls-from-curvature-volume-bounds-existence-noncollapsing}
\source{kleiner-lott}
\uses{prop:kl-proof-double-sided-curvature-bound-thick-part-modulo-two-propositions-existence-canonical-neighborhoods, rem:kl-no-local-collapsing-propagates-forward-time-larger-scales-estimates-no-local-collapsing-propagates-forward-time-larger-scales, thm:kl-no-local-collapsing-propagates-forward-time-larger-scales-existence-kappa-solutions}
\dcref{Remark 53.2}
\group{Kleiner--Lott Canonical Neighborhoods}

\label{rmkI.12.2}

The restriction $r_0^2\Phi(r_0^{-2})<\rho$ is necessary for the conclusion;
the volume hypothesis supplies noncollapsing.
\end{remark}
\begin{lemma}[existence for kappa solutions]
\label{lem:kl-later-scalar-curvature-bounds-bigger-balls-from-curvature-volume-bounds-existence-kappa-solutions}
\source{kleiner-lott}
\dcref{Lemma 53.3}
\group{Kleiner--Lott Canonical Neighborhoods}

\label{12.2lemma}
For every $\epsilon>0$ there is $K=K(A,\epsilon)<\infty$ such that every
$x$ with $d_{r_0^2}(x,x_0)<Ar_0$ satisfies either
$R(x,r_0^2)<Kr_0^{-2}$, or, with $Q=R(x,r_0^2)$, the region
\[
\{(x',t'):d_t(x',x)<(\epsilon Q)^{-1},\
t-(\epsilon Q)^{-1}\leq t'\leq t\},\quad t=r_0^2,
\]
after scaling by $Q$, is $\epsilon$-close to a $\kappa$-solution region.
\end{lemma}
\begin{remark}[structural remark on thick--thin geometry]
\label{rem:kl-later-scalar-curvature-bounds-bigger-balls-from-curvature-volume-bounds-structural-remark-thick-thin-geometry}
\source{kleiner-lott}
\uses{thm:kl-canonical-neighborhood-theorem-compactness-canonical-neighborhoods}
\dcref{Remark 53.4}
\group{Kleiner--Lott Canonical Neighborhoods}


This lemma is the localized form of Theorem \ref{thmI.12.1}: both its
hypotheses and its conclusion are centered at $x_0$.
\end{remark}
\begin{proof}
\proves{lem:kl-later-scalar-curvature-bounds-bigger-balls-from-curvature-volume-bounds-existence-kappa-solutions}
\uses{lem:kl-canonical-neighborhood-theorem-canonical-neighborhoods-lemma-9,lem:kl-length-distortion-estimates-distance-distortion-lemma,thm:kl-canonical-neighborhood-theorem-compactness-canonical-neighborhoods,thm:kl-no-local-collapsing-propagates-forward-time-larger-scales-existence-kappa-solutions}
\sketch Point-pick a bad sequence so nearby points of twice the curvature are
good.  Lemma \ref{rescaled} yields local curvature bounds and a pointed
limit; extend it backward using distance distortion and noncollapsing.  The
resulting ancient convergence contradicts the selected point's failure of
the second alternative.
\end{proof}
\subsection{Earlier scalar curvature bounds on smaller balls from lower curvature bounds and volume bounds}
\dcref{I.12.3}
\label{I.12.3}
\begin{lemma}[existence for scalar curvature]
\label{lem:kl-earlier-scalar-curvature-bounds-smaller-balls-from-lower-curvature-bound-existence-scalar-curvature}
\source{kleiner-lott}
\dcref{Lemma 54.1}
\group{Kleiner--Lott Canonical Neighborhoods}

\label{Alexlemma}
Given $w'>0$ and $n\in\mathbb Z^+$, there is $c=c(w',n)>0$ such that a
compactly contained radius-$r$ ball with sectional curvature at least
$-r^{-2}$ and volume at least $w'r^n$ contains a ball of radius
$r'\geq cr$ and volume at least $\frac12\omega_n(r')^n$.
\end{lemma}
\begin{proof}
\sketch Rescale to $r=1$ and take a pointed Gromov--Hausdorff limit.
Weak convergence of volume measures gives a small ball with almost Euclidean
volume, contradicting the assumed failure.
\end{proof}
\begin{theorem}[estimates for scalar curvature]
\label{thm:kl-earlier-scalar-curvature-bounds-smaller-balls-from-lower-curvature-bound-estimates-scalar-curvature}
\source{kleiner-lott}
\dcref{Theorem 54.2}
\group{Kleiner--Lott Canonical Neighborhoods}

(cf. Theorem I.12.3) \label{thmI.12.3}
For every $w>0$ there are $\tau(w)>0$, $K(w)<\infty$, and $\rho(w)>0$ such
that, for a $\Phi$-almost nonnegatively curved flow on a closed
three-manifold, $t_0\geq4\tau r_0^2$,
$r_0^2\Phi(r_0^{-2})<\rho$, volume at least $wr_0^3$, and sectional
curvature at least $-r_0^{-2}$ on the time-$t_0$ ball, one has
\[
R(x,t)\leq Kr_0^{-2}
\]
for $t\in[t_0-\tau r_0^2,t_0]$ and $d_t(x,x_0)\leq r_0/4$.
\end{theorem}
\begin{proof}
\uses{cor:kl-curvature-bounds-ricci-flow-solutions-nonnegative-curvature-operator-ass-existence-curvature-operators-3,lem:kl-earlier-scalar-curvature-bounds-smaller-balls-from-lower-curvature-bound-existence-scalar-curvature,rem:kl-later-scalar-curvature-bounds-bigger-balls-from-curvature-volume-bounds-existence-noncollapsing,thm:kl-later-scalar-curvature-bounds-bigger-balls-from-curvature-volume-bounds-positivity-scalar-curvature}
\sketch Set $\tau=\tau_0/2$ and $K=C+2B/\tau_0$.  Point-pick a
counterexample.  If the lower curvature bound persists for $2\tau r_0^2$,
Corollary \ref{I.11.6end} gives the result.  At first failure, distortion and
\ref{Alexlemma} produce an almost-Euclidean subball of controlled radius;
Theorem \ref{thmI.12.2} improves its curvature and pinching restores a lower
bound better than $-r_0^{-2}$, contradicting maximality.
\end{proof}

\subsection{Small balls with strongly negative curvature are volume-collapsed}
\dcref{I.12.4}
\label{I.12.4}
\begin{corollary}[positivity for curvature]
\label{cor:kl-small-balls-strongly-negative-curvature-volume-collapsed-positivity-curvature}
\source{kleiner-lott}
\dcref{Corollary 55.1}
\group{Kleiner--Lott Canonical Neighborhoods}

(cf. Corollary I.12.4) \label{corI.12.4}
For every $w>0$ there is $\rho>0$ such that a $\Phi$-almost nonnegatively
curved flow on a closed three-manifold, defined on $[0,T)$ with $T\geq1$,
satisfies
\[
\inf_{B_{t_0}(x_0,r_0)}\Rm=-r_0^{-2}
\quad\Longrightarrow\quad
\vol(B_{t_0}(x_0,r_0))\leq wr_0^3
\]
whenever $t_0\geq1$ and $r_0<\rho$.
\end{corollary}
\begin{proof}
\uses{thm:kl-earlier-scalar-curvature-bounds-smaller-balls-from-lower-curvature-bound-estimates-scalar-curvature,thm:kl-later-scalar-curvature-bounds-bigger-balls-from-curvature-volume-bounds-positivity-scalar-curvature}
\sketch A larger volume would give, by Theorem \ref{thmI.12.3}, a
two-sided bound on a smaller earlier ball.  Distortion and Theorem
\ref{thmI.12.2} then give an upper scalar bound on the original ball, while
\eqref{phieqnn} gives $\Rm> -r_0^{-2}$ for sufficiently small $r_0$, a
contradiction.
\end{proof}

\section{Thick-thin decomposition for nonsingular flows} \label{I.13.1}
\dcref{I.13.1}
For a nonflat immortal three-dimensional flow, let $\widehat r(x,t)$ be the
unique radius with
\[
\inf_{B_t(x,\widehat r(x,t))}\Rm=-\widehat r(x,t)^{-2},
\]
and define
\begin{equation}\label{thinpart}
M_{\thin}(w,t)=\{x:\vol(B_t(x,\widehat r(x,t)))<
w\widehat r(x,t)^3\},\qquad M_{\thick}(w,t)=M-M_{\thin}(w,t).
\end{equation}
For large $t$, Corollary \ref{corI.12.4} implies
$x\in M_{\thick}(w,t)\Rightarrow\widehat r(x,t)\geq\widehat\rho(w)\sqrt t$.
\begin{theorem}[existence for thick--thin geometry]
\label{thm:kl-thick-thin-decomposition-nonsingular-flows-existence-thick-thin-geometry}
\source{kleiner-lott}
\dcref{Theorem 56.2}
\group{Kleiner--Lott Thick--Thin Decomposition}

(cf. I.13.1) \label{thmI.13.1}
For every $w>0$ there are $T=T(w)>0$, $\overline\rho=\overline\rho(w)>0$,
and $K=K(w)<\infty$ such that, if $t\geq T$ and
$x\in M_{\thick}(w,t)$, then
\[
|\Rm|\leq Kt^{-1}\quad\hbox{on }B_t(x,\overline\rho\sqrt t),
\qquad
\vol(B_t(x,\overline\rho\sqrt t))
\geq\frac1{10}w(\overline\rho\sqrt t)^3.
\]
\end{theorem}
\begin{proof}
\uses{cor:kl-small-balls-strongly-negative-curvature-volume-collapsed-positivity-curvature,thm:kl-earlier-scalar-curvature-bounds-smaller-balls-from-lower-curvature-bound-estimates-scalar-curvature,thm:kl-later-scalar-curvature-bounds-bigger-balls-from-curvature-volume-bounds-positivity-scalar-curvature}
\sketch Put $r=c\widehat\rho\sqrt t$.  The thick hypothesis and
Bishop--Gromov give
\[
\vol(B_t(x,r))\geq
\frac{w(c\widehat\rho)^3t^{3/2}}{3\int_0^1\sinh^2u\,du}
\geq\frac1{10}w(c\widehat\rho)^3t^{3/2}.
\]
For $c$ small and $t$ large, Theorem \ref{thmI.12.3}, pinching, distortion,
and Theorem \ref{thmI.12.2} give $|\Rm|\leq K(w)t^{-1}$ on the controlled
ball and preserve the displayed volume lower bound.  Set
$\overline\rho=c\widehat\rho$.
\end{proof}
\begin{remark}[estimates for thick--thin geometry]
\label{rem:kl-thick-thin-decomposition-nonsingular-flows-estimates-thick-thin-geometry}
\source{kleiner-lott}
\uses{thm:kl-later-scalar-curvature-bounds-bigger-balls-from-curvature-volume-bounds-positivity-scalar-curvature, thm:kl-thick-thin-decomposition-nonsingular-flows-existence-thick-thin-geometry}
\dcref{Remark 56.5}
\group{Kleiner--Lott Thick--Thin Decomposition}


For every $A>0$, the same argument gives
$|\Rm|(x,t)\leq K(A,w)t^{-1}$ on
$B_t(x,A\overline\rho\sqrt t)$ for thick points and large $t$.
\end{remark}
